% !TeX spellcheck = nb_NO
% Kommentaren ovenfor gir instruks til editoren din om at den skal bruke norsk stavekontroll. Dette fungerer bl.a. i TeXstudio.

\documentclass[a4paper,norsk]{article}
% Norsk orddeling
\usepackage[norsk]{babel}
% AMS-pakkene er nyttige
\usepackage{amsmath,amsfonts,amssymb,amsthm}
% For å definere lister, som f.eks. deloppgaver
\usepackage{enumitem}
% Inneholder mye nyttig, slik som \coloneqq
\usepackage{mathtools}
% Penere kalligrafiske fonter
\usepackage[utf8]{inputenc}
\usepackage[T1]{fontenc}
\usepackage[a4paper, margin=0.5in]{geometry}
\usepackage{mathabx}
\usepackage{listings}
\usepackage{xcolor}
\usepackage{ifthen}
\usepackage{parskip}

\usepackage[most]{tcolorbox}
\tcbuselibrary{theorems, breakable, skins}
\tcbset{before upper=\par, after upper=\par} % allows itemize and enumerate inside boxes

\usepackage{hyperref}

\setlength{\parindent}{0pt}
\setlength{\parskip}{6pt plus 2pt minus 1pt}


\tcbset{
  % Better page breaking for all tcolorbox environments
  enhanced,
  breakable,
  before skip=5pt plus 2pt minus 1pt,
  after skip=5pt plus 2pt minus 1pt,
  parbox=false,
}
\newtcbtheorem[]{definitionbox}{Definisjon}{
  colback=blue!3,
  colframe=blue!70!black,
  fonttitle=\bfseries\color{white},
  coltitle=white,
  boxrule=0.6pt,
  arc=2mm,
  enhanced,
  breakable
}{def}

\newtcbtheorem[]{theorembox}{Teorem}{
  colback=green!3,
  colframe=green!50!black,
  fonttitle=\bfseries,
  coltitle=black,
  boxrule=0.6pt,
  before upper=\par\ignorespaces,
  after upper=\par\unskip,
  arc=2mm,
  enhanced,
  breakable
}{thm}

\newtcbtheorem[]{oppgavebox}{Oppgave}{
  colback=gray!5,
  colframe=black!60,
  fonttitle=\bfseries,
  coltitle=black,
  boxrule=0.6pt,
  arc=2mm,
  enhanced,
  breakable,
  parbox=false
}{oppg}

\newtcolorbox{sectionbox}[1][]{
  colback=white,
  colframe=black!70,
  boxrule=0.7pt,
  arc=3mm,
  top=2mm, bottom=2mm, left=2mm, right=2mm,
  title=#1,
  enhanced,
  breakable
}

\newtcbtheorem[]{corollarybox}{Korollar}{
  colback=purple!5,
  colframe=purple!60!black,
  fonttitle=\bfseries\color{white},
  coltitle=white,
  boxrule=0.6pt,
  arc=2mm,
  enhanced,
  breakable,
  before upper=\par\ignorespaces,
  after upper=\par\unskip
}{cor}

\newtcbtheorem[]{lemmabox}{Lemma}{
  colback=orange!5,
  colframe=orange!70!black,
  fonttitle=\bfseries,
  boxrule=0.6pt,
  arc=2mm,
  enhanced,
  breakable
}{lem}

\newtcbtheorem[]{propositionbox}{Proposisjon}{
  colback=cyan!5,
  colframe=cyan!60!black,
  fonttitle=\bfseries,
  coltitle=black,
  boxrule=0.6pt,
  arc=2mm,
  enhanced,
  breakable
}{prop}

\newtcbtheorem[]{proofbox}{Bevis}{
  enhanced,
  breakable,
  colback=white,
  colframe=black,
  fonttitle=\bfseries\color{white},
  coltitle=white,
  colbacktitle=black,
  title style={opacity=0.9},
  boxrule=0.6pt,
  arc=2mm,
  before upper=\par\ignorespaces,
  after upper=\par\unskip,
  parbox=false, % Viktig for sidebryting
  segmentation style={black,solid,opacity=0.2,line width=0.5pt}
}{proof}

\newlist{suboppgave}{enumerate}{1}
\setlist[suboppgave,1]{
  label=\textbf{(\alph*)},
  ref=\theenumi\alph*,
  leftmargin=1.5em,
  itemsep=0.4em,
  topsep=0.3em,
  parsep=0pt,
  partopsep=0pt
}



\setcounter{secnumdepth}{2}  % Sections and subsections numbered

% Python style  
\lstdefinestyle{pythonstyle}{
    language=Python,
    basicstyle=\ttfamily\small,
    keywordstyle=\color{blue},
    commentstyle=\color{green},
    numbers=left,
    frame=single
}

% Bruk bedre symboler for ≥, ≤, epsilon og phi
\renewcommand{\geq}{\geqslant}
\renewcommand{\leq}{\leqslant}
\newcommand{\eps} {\varepsilon}
\renewcommand{\epsilon}{\varepsilon}
\renewcommand{\phi}{\varphi}

% Definer de vanligste mengdene og funksjonene
\newcommand{\R}{\mathbb{R}}
\newcommand{\K}{\mathbb{K}}
\newcommand{\C}{\mathbb{C}}
\newcommand{\N}{\mathbb{N}}
\newcommand{\Z}{\mathbb{Z}}
\newcommand{\Q}{\mathbb{Q}}
% \L er definert som noe annet fra før av, så vi må redefinere \L
\renewcommand{\L}{\mathcal{L}}
% Rommet av polynomer
\newcommand{\Pol}{{\mathcal{P}}}
\DeclareMathOperator{\id}{id}
\DeclareMathOperator{\Image}{im}
\DeclareMathOperator{\Kernel}{ker}
\DeclareMathOperator{\Span}{span}
\DeclareMathOperator{\Rank}{rank}
\DeclareMathOperator{\Diag}{diag}
\DeclareMathOperator{\Proj}{proj}
\newcommand{\basis}{{\mathcal{B}}}
\newcommand{\basiss}{{\mathcal{C}}}
\newcommand{\basisss}{{\mathcal{D}}}

\newtcolorbox{algoritme}[1]{
  title=Algoritme~#1,
  colback=white,
  colframe=black,
  boxrule=0.5pt,
  arc=2mm,
  left=3mm,
  right=3mm,
  top=1mm,
  bottom=1mm
}

% Definer oppgavemiljøet
\newenvironment{oppgave}[1]%
	{\trivlist{}\item\relax\textbf{Løsning på #1.}\medspace}%
	{\endtrivlist}
% Deloppgavelister
\newlist{deloppgave}{enumerate}{1}
\setlist[deloppgave,1]{label=(\textbf{\alph*}),align=left}

\title{Real Analysis \\ MAT2400 University of Oslo, Spring 2026}
\author{Oliver Ekeberg}

\begin{document}
\tableofcontents


\section*{3.1 Definitions and examples}
\addcontentsline{toc}{section}{Defs and examples}


\begin{definitionbox}{3.1.1}{}
A metric space $ \left( X, d \right)  $ consists of a nonempty set X and a function 
\[
    d: X \times X \rightarrow \left[ 0, \infty \right)
\]
such that 
\begin{itemize}
    \item $ \forall x, y \in X $ we have $ d(x, y) \geq 0 $ and $ d(x,y) = 0 \iff x=y $ 
    \item $ \forall x, y \in X $ we have $ d(x,y) = d(y,x) $ 
    \item $ \forall x, y, z \in X $ we have $ d(x,y) \leq d(x,z) + d(z,y) $ 
\end{itemize}
A function d statisfying these conditions is called a metric on X
\end{definitionbox}

\begin{oppgavebox}{3.1.1}{}
We shall now take a look at an example of a different kind. Assume that we want to send messages in a language with N symbols (letters, numbers, punctuation marks, space, etc.). We assume that all messages have the same length K (if they are too short or too long, we either fill them out or break them into pieces). We let X be the set of all messages, i.e., all sequences of symbols from the language of length K. If x = (x1,x2,...,xK) and y = (y1,y2,...,yK) are two messages, we define
$$d(x, y) = \text{the number of indices n such that } x_n = y_n. $$
It is not hard to check that d is a metric. It is usually referred to as the Hamming- metric, and is much used in communication theory, where it serves as a measure of how much a message gets distorted during transmission.

\begin{proofbox}{3.1.1}{}
    This set X can be called a metric space if it satisfies the three conditions of def 3.1.1. We have that the distance between x and y is the number of indices n such that $ x_n \neq y_n $. We have that d satisfies the first condition if and only if $ x_n = y_n $ for all indeces n. This means that $ x = y $ where
    \[
        x = \left( x_1, x_2, ..., x_K \right) 
        y = \left( y_1, y_2, ..., y_K \right) 
    \]
    And this is equivalent with $ d(x,y) = 0 $. Since $ x=y \Rightarrow d(x,y) = d(x,x) = 0 $ \\

    We now prove the second condition. $ \forall x, y \in X $, $ d(x, y) = d(y, x) $.
    
    If $ d =  $ \# of indices s.t. $ x_n \neq y_n $, we have that 
    \[
        d(x, y) = \text{\# of indices such that } x_n \neq y_n
    \]
    which is equivalent with 
    \[
        d(y, x) = \text{\# of indices such that } y_n \neq x_n
    \]
    So $ d(x, y) = d(y, x) $\\

    Now for the third condition. $ \forall x, y, z \in X $ we have $ d(x, y) \leq d(x, z) + d(z, y) $.
    \begin{align*}
      x &= \left( x_1, x_2, ..., x_K \right) \\
      y &= \left( y_1, y_2, ..., y_K \right) \\
      z &= \left( z_1, z_2, ..., z_n \right) 
    \end{align*}
    Then 
    \begin{align*}
      d(x, y) &= \text{\# of indices s.t. } x_n \neq y_n\\
      d(x, z) = \text{\# of indices s.t. } x_n \neq z_n
      d(z, y) = \text{\# of indices s.t. } z_n \neq y_n
    \end{align*}
    We know that if $ d(x, z) = k, d(z, y) = r $ and $ g = d(x, y) $, then $ r + k \geq g $ and $ r+k=g \quad \iff \quad x=y $. 
    
\end{proofbox}
\end{oppgavebox}


\begin{oppgavebox}{1.3.2}{}
Prove that $ \left( \mathbb{R}^n, d \right)  $ is a metric space, with 
\[
    d(x, y) = \left\lVert x - y \right\rVert = \sqrt{ \sum_{ n }^{ i=1 } \left( x_i - y_i \right) ^2  }
\]


\begin{proofbox}{1.3.2}{}
  \begin{itemize}
    \item The first condition is that $ d(x, y) \geq 0 $ and that $ d(x, y) = 0 \iff x = y $. We know that $ d(x, y)  $  is obviously larger than 0, and 
    \[
        d(x, x) = \sqrt{ \sum_{ n }^{ i=1 } \left( x_i - x_i \right)^2  } = \sqrt{ \sum_{ n }^{ i=1 } 0  } = 0
    \]

    \item The second condition is that $ d(x, y) = d(y, x) $. 
    \begin{align*}
      d(x, y) &= \sqrt{ \sum_{ n }^{ i=1 } \left( x_i - y_i \right) ^2 } \\
      &= \sqrt{ \sum_{ n }^{ i=1 } \left( -1 \right)^2  \left( y_i - x_i \right) ^2 } \\
      &= d(x, y)
    \end{align*}
     
    \item For the third condition, we have that 
    \begin{align*}
      d(x, y) &= \left\lVert x - y \right\rVert \\
      &= \left\lVert \left( x - z \right) + \left( z - y \right)  \right\rVert \leq \left\lVert x-z \right\rVert + \left\lVert z-y \right\rVert \\
      &= d(x,z) + d(z,y)
    \end{align*}
    
  \end{itemize}
\end{proofbox}
\end{oppgavebox}

\begin{propositionbox}{3.1.4}{Inverse Triangle inequality}
$ \forall x, y, z \in X $ with $ \left( X, d \right)  $ metric space, 
\[
    \left| d(x, y) - d(x,z) \right| \leq d(y, z)
\]

\begin{proofbox}{3.1.4}{}
  We have that $ d(x, y) \leq d(x, z) + d(z, y) $. Then $ d(x, y) - d(x, z) \leq d(y, z) $. Now we need to prove the negative.

  \[
      d(x, z) \leq d(x, y) + d(y, z) \Rightarrow d(x, z) - d(x, y) \leq d(y, z)
  \]
  and thus
  \[
      d(x, y) - d(x, z) = d(x, z) - d(x, y) \leq d(y, z) \Rightarrow \left| d(x, y) - d(x, z)  \right| \leq d(y, z)
  \]
\end{proofbox}
\end{propositionbox}

\begin{oppgavebox}{3.1.6}{}
If V is a vector space over R or C, a function $ \left\lVert \cdot \right\rVert : V \rightarrow \mathbb{R} $  is called a norm if the following conditions are satisfied:
\begin{itemize}
  \item (i) $ \forall x \in V, \left\lVert x \right\rVert \geq 0$ and $ \left\lVert x \right\rVert = 0 \iff x=0 $ 
  \item $ \left\lVert \alpha_{}^{} x \right\rVert  = \left| \alpha_{}^{}  \right| \left\lVert x \right\rVert \quad \forall x \in V $ and $ \alpha_{}^{} \in \mathbb{K} $ 
  \item $ \left\lVert x + y \right\rVert \leq \left\lVert x \right\rVert + \left\lVert y \right\rVert \quad \forall x, y \in V $  
\end{itemize}


\begin{proofbox}{3.1.6}{}
  \begin{itemize}
    \item Since $ \left\lVert \cdot \right\rVert : V \rightarrow \mathbb{R}  $, then it satisfies all three conditions for a norm. To show that it satisfies the first condition of a metric, then

  \[
    \left\lVert x \right\rVert = d(x, 0) = \left\lVert x - 0 \right\rVert \geq 0
  \]
  and this becomes similar to the prove for norm

  \item Have to now show that $ d(x, y) = d(y, x) $ 
  \begin{align*}
    d(x, y) &= \left\lVert x - y \right\rVert \\
    &= \left\lVert \left( -1 \right) \left( y - x \right)  \right\rVert \\
    &= \left| \left( -1 \right) \right| \left\lVert y - x \right\rVert \\
    &= \left\lVert y - x \right\rVert = d(y, x)
  \end{align*}

  \item Have to now show that $ d(x, y) \leq d(x, z) + d(z, y) $ 
  \begin{align*}
    d(x, y) &= \left\lVert x - y \right\rVert \\
    \left\lVert \left( x - z \right) + \left( z-y \right)  \right\rVert \leq \left\lVert x-z \right\rVert + \left\lVert z-y \right\rVert \\
    &= d(x, z) + d(z, y)
  \end{align*}
  
  And the proof is finished
  \end{itemize} 
  
\end{proofbox}
\end{oppgavebox}

\begin{oppgavebox}{3.1.7}{}
Show that if $ \left( x_1, x_2, ..., x_n \right)  $ are points in a metric space, then 
\[
    d(x_1, x_n) \leq d(x_1, x_2) + d(x_2, x_3) + ... + d(x_{n-1} , x_n)
\]

\begin{proofbox}{3.1.7}{}
  We use induction to prove this. Assume that 
  \[
      d(x_k, x_{k+2}) \leq d(x_k, x_{k+1}) + d(x_{k+1}, x_{k+1})
  \]
  Then for $ k=0 $, we get that 
  \[
      d(x_1, x_3) \leq d(x_1, x_2) + d(x_2, x_3)
  \]
  
  For $ k=k $ we assume that 
  \[
      d(x_k, x_{k+2}) \leq d(x_k, x_{k+1}) + d(x_{k+1}, x_{k+1})
  \]
  and now we apply our hypothesis, that 
  \[
      d(x_k, x_{k+3}) \leq d(x_k, x_{k+2}) + d(x_{k+2}, x_{k+3})
  \]
  since $ d(x_k, x_{k+2}) \leq d(x_k, x_{k+1}) + d(x_{k+1}, x_{k+1}) $, we do get that 
  
  \[
      d(x_k, x_{k+3}) \leq d(x_k, x_{k+2}) + d(x_{k+2}, x_{k+3})
  \]
  and this will now apply for all k. Thus the statement that 
  \[
      d(x_1, x_n) \leq d(x_1, x_2) + ... + d(x_{n-1}, x_n)
  \]
  is indeed true  
\end{proofbox}
\end{oppgavebox}


\section*{3.2 Convergence and continuity}
\addcontentsline{toc}{section}{Convergence and continuity}

\begin{definitionbox}{3.2.1}{}
Let $ \left( X, d \right)  $ be a metric space. A sequence $ \left\{ x_n \right\} _{ n \in \mathbb{N} } $ converges to a point $ a \in X $ if there for every $ \epsilon > 0 $ exists an $ N \in \mathbb{N} $ such that $ d(x_n, a) < \epsilon \forall n \geq N $ 
\end{definitionbox}

\begin{lemmabox}{3.2.2}{}
A sequence $ \left\{ x_n \right\} _{n \in \mathbb{N} } $ in a metric space $ \left( X, d \right)  $ converges to a point a, if and only if $ d(x_n, 0)  $ converges to zero

\begin{proofbox}{3.2.2}{}
  The sequence $ d( \left\{ x_n \right\} , a) $ converges to a point $ 0 \in X $ if there is for every $ \epsilon > 0 \quad \exists N \in  \mathbb{N} $ such that $ \forall n \geq N $, $ d(x_n, a) < \epsilon  $  
\end{proofbox}
\end{lemmabox}

Can a sequence converge towards multiple points? 

\begin{propositionbox}{3.2.3}{}
A sequence in a metric cannot converge to more than one point.

\begin{proofbox}{3.2.3}{}
  Assume that $ \lim_{n \to \infty} x_n = a  $ and $ \lim_{n \to \infty} x_n = b $. If we take the triangle inequality on $ d(a, b) $, we get
  
  \[
    d(a, b) \leq d(a, x_n) + d(x_n, b)
  \]
  We know that the sequences $ d(a, x_n) = d(x_n, a) $ and $ d(x_n, b) $  converges to zero if $ x_n $ converges to a and b respectively. This

  \begin{align*}
    d(a, b) \leq \lim_{n \to \infty} d(x_n, a) + d(x_n, b) \\
    &= 0+0
  \end{align*}
  And surely, $ d(a, b) $ cannot be less than zero. Thus, $ d(a, b) = 0 \iff a=b $
\end{proofbox}
\end{propositionbox}



\begin{definitionbox}{3.2.4}{}
Assume $ \left( X, d_X \right)  $ and $ \left( Y, d_Y \right)  $ are two metric spaces. A function $ f: X \rightarrow Y $ is continuous at a point $ a \in X $ if for every $ \epsilon > 0 \exists \delta > 0 $ s.t. $ d_X (x, a) < \delta \Rightarrow d_Y (f(x), f(a)) < \epsilon $  
\end{definitionbox}

\begin{propositionbox}{3.2.5}{}
The following is equivalent for a function $ f: X \rightarrow Y $ for two metric spaces.

\begin{itemize}
  \item f is continuous at a point $ a \in X $ 
  \item For all sequences $ \left\{ x_n \right\}  $ converging to a, the sequences $ \left\{ f(x_n) \right\}  $ are converging to $ f(a) $  
\end{itemize}

\begin{proofbox}{3.2.5}{}
  First we show that $ i) \Rightarrow ii) $. Assume that f is continuous at a. We want to show that 
  \[
      \forall \delta > 0, \exists N \in \mathbb{N} \text{ such that } \forall n\geq N, d_X(x_n, a) < \delta
  \]
   
  We know that $ \forall \epsilon > 0 \exists \delta > 0  $ s.t. $ d_X \left( x_n, a \right) < \delta \Rightarrow d_Y \left( f(x_n), f(a) \right) < \epsilon   $. Since $ x_n \rightarrow a $, we know that $ \exists N \in \mathbb{N} $ such that $ \forall n \geq N, d_X(x_n, a) < \delta $. But then $ d_Y (f(x_n), f(a)) < \epsilon  $. So $ f:X \rightarrow Y $ preserves convergence\\

  Second we want to prove that $ \neg i) \Rightarrow ii) $ to show a contradiction. Assume that f is \textit{not} continuous at a. Want to show that for all converging sequences $ \left\{ x_n \right\}  $ to a, in the metric space X, the sequence $ \left\{ f(x_n) \right\}  $ does not converge to $ f(a) $ in the metric space Y. 

  We know that $ x_n $ converges to a. This means that for evey $ \epsilon >0 \exists N \in \mathbb{N} $ such that $ d_X(x_n, a) < \epsilon. \forall n \geq N $. Since f is not cont at a, 
  
  \[
      \exists \epsilon > 0 \text{ such that } \forall \delta > 0, \exists x \in X \text{ such that } d_X(x_n, a) < \delta \Rightarrow d_Y(f(x_n), f(a)) \geq \epsilon   
  \]

  Because of the fact that f is not continuous at a, we can pick $ \delta = \epsilon  $ and get that $ \forall \epsilon > 0, \exists N \in \mathbb{N} $ such that $ d_Y(f(x_n), f(a)) \geq \epsilon  \quad \forall n \geq N$ 
  
\end{proofbox}
\end{propositionbox}

\begin{oppgavebox}{3.2.1}{}
Assume that $(X, d)$ is a discrete metric space (recall Example 6 in Section 3.1). Show that the sequence $ \left\{ x_n \right\} $ converges to a if and only if there is an $ N \in \mathbb{N} $  such that $ x_n = a \forall n \geq N $. 

\begin{proofbox}{3.2.1}{}
  We have two events 
  \begin{itemize}
    \item A: THe sequence $ \left\{ x_n \right\}  $ converges to a. 
    \item B: There is an $ N \in \mathbb{N} $ such that $ x_n = a\forall n \geq N $ 
  \end{itemize}
  Our task is to prove $ A \iff B $. We begin with $ A \Rightarrow B $  Assume first that the sequcence $ \left\{ x_n \right\}  $ converges to a. This means that for every $ \epsilon > 0 $ there exists some $ N \in \mathbb{N} $ such that $ \forall n \geq N $, we have that 
  \[
      d(x_n, a) < \epsilon 
  \]

  where $ d(x_n, a) = \begin{cases}
    0, x_n =a \\
    1, x_n \neq a
  \end{cases} $ 
  We can pick $ 0 < \epsilon < 1 $ and get that $ \forall n \geq N $, $ d(x_n, a) < 1 $, but then $ d(x_n, a) = 0 $ due to the discrete metric. And since a metric is equal to 0 if and only if $ x_n = a $, then there is an $ N \in  \mathbb{N} $ such that $ \forall n \geq N $, we have that $ x_n = a $         \\

  We now prove $ B \Rightarrow A $. Since $ \exists N \in  \mathbb{N} $ such that $ x_n = a \quad \forall n \geq N $, then $ \forall \epsilon > 0 $, we have that  $ d(x_n, a) = 0 \quad < \epsilon  \forall n \geq N$, and therefore, the sequence $ \left\{ x_n \right\}  $ converges to a
\end{proofbox}
\end{oppgavebox}

\begin{oppgavebox}{3.2.2}{}
Prove proposition 3.2.6 without using proposition 3.2.5.

\begin{propositionbox}{3.2.6}{}
Let $ \left( X, d_X \right), \left( Y, d_Y \right) , \left( \mathbb{Z}, d_Z \right)   $ be three metric spaces. Assume that 
\begin{itemize}
  \item $ f: X  \rightarrow Y $ 
  \item $ g : Y  \rightarrow Z $ 
  \item $ h : X  \rightarrow Z $ 
\end{itemize}
be the composition $ h(x) = g \left( f(x) \right)  $. SHow that if f is continious at $ a \in X $, and that g is continious at  $ b = f(a) $, then h is continuous at a

\begin{proofbox}{3.2.6}{}
  If f is continous at $ a \in X $, then $ \forall \epsilon > 0 \quad \exists \delta > 0 s.t. \forall x \in X, d_X(x, a) < \delta \Rightarrow d_Y( f(x), f(a) ) < \epsilon  $. If g is continuous at $ b = f(a) \in Y $, then $ \forall \epsilon > 0 \quad \exists \delta > 0 s.t. \forall y \in Y, d_Y(y, b) < \delta \Rightarrow d_Z(g(y), g(b)) < \epsilon  $. \\

  This means that $ h :=  g \circ f : X \Rightarrow Z $ is continuous if $ \forall \epsilon > 0 \exists \delta > 0 s.t. \forall x \in X, d_X(x, a) < \delta \Rightarrow d_Z (h(x), h(a)) < \epsilon  $ where $ h(a) = g(f(a)) = g(b) $ 
\end{proofbox}
\end{propositionbox}
\end{oppgavebox}

\begin{oppgavebox}{3.2.5}{}
Let $ \left( X, d \right)  $ be a metric space and choose $ a \in X $. Show that $ f: X \Rightarrow \mathbb{R} $ given by $ f(x) = d(x, a) $ is continuous. Use the metric $ d_{ \mathbb{R}} (x, y) = \left| x - y \right|  $ 

\begin{proofbox}{3.2.5}{}
  f is continious at $ a \in X $ if $ \forall \epsilon  > 0 \exists \delta > 0 s.t. \forall x \in  X, d_X(x,a) < \delta \Rightarrow d_{ \mathbb{R}}(f(x), f(a)) < \epsilon  $. 
  \begin{align*}
    d(x, a) &= \left| x - a \right| 
  \end{align*}

  \begin{align*}
    d( f(x), f(a)) &= \left| \left| x-a \right|  - \left| a-a \right|  \right| \\
    &= \left| x-a \right| 
  \end{align*}
  So we get that f is continuous at a if $ \forall \epsilon > 0 \exists \delta > 0 s.t. \forall x \in X, \left| x-a \right| < \delta \Rightarrow \left| x-a \right| < \epsilon  $. Just pick $ \delta = \epsilon  $, and the proof is complete.
  
  
\end{proofbox}
\end{oppgavebox}

\begin{oppgavebox}{3.2.6}{}
Let $ \left( X, d_X \right) , \left( Y, d_Y \right)  $ be two metric spaces. A function $ f: X \rightarrow Y $ is said to be a Lipschitz function if $ \exists K \in \mathbb{K}  $ such that 
\[
    d_Y (f(u), f(v)) \leq K d_X (u, v) \quad \forall u, v \in X
\]
Show that all Lipschitz functions are continuous
\begin{proofbox}{3.2.6}{}
  We have that 
  \begin{itemize}
    \item A = f is Lipschitz
    \item B = f is continuous
  \end{itemize}
  We are tasked to prove $ A \Rightarrow B $. Assume that $ f: X \rightarrow Y $ is Lipschitz. This means that $ \exists K \in \mathbb{R}, K > 0 $ such that $ d_Y (f(u), f(v)) \leq K d_X( y, v) \forall u, v \in X $. A function is continuous at $ v \in X $  if $ \forall \epsilon > 0 \exists \delta > 0 s.t. \forall u \in X, d_X(u, v) < \delta \Rightarrow d_Y (f(u), f(v)) < \epsilon   $. We can pick $ \delta = \frac{ \epsilon }{ K } $ and get that 
  \[
      \forall \epsilon > 0 \exists \delta > 0 s.t. \forall u \in X, d_X (u, v) < \delta = \frac{ \epsilon  }{ K } \Rightarrow d_Y (f(u), f(v)) \leq K d_X(u, v) < K \frac{ \epsilon }{ K } = \epsilon 
  \]
  Therefore, f is continous at v.
\end{proofbox}

\end{oppgavebox}

\begin{oppgavebox}{3.2.8}{}
\begin{itemize}
  \item a) Assume that $ \left\{ x_n \right\}  $ is a sequence in X converging to x. Show that $ d(x_n, y) \rightarrow d(x, y) $
  \begin{proofbox}{3.2.8 a)}{}
    Since $ \left\{ x_n \right\}  $ is a sequence converging to $ x \in X $, it is equivalent with saying that $ d(x, x_n) \rightarrow 0 $. This is because if $ \left\{ x_n \right\}  $ is a sequence converging to $ x \in X $, then 
    \[
        \forall \epsilon > 0, \exists N \in \mathbb{N} \text{ such that } \forall n \geq N, \text{da er } d(x_n, x) < \epsilon 
    \]
    And since we can make $ \epsilon  $ arbitrarily small, we say that the sequence $ \left\{ d \left( x_n, x \right)  \right\}  $ converges to $ 0 $ . To show equivalence the other way, assume that $ \left\{ d \left( x_n, x \right)  \right\}  $ converges to $ 0 $. Then 
    \[
        \forall \epsilon > 0, \exists N \in \mathbb{N} \text{ such that } \forall n \geq N, \left| d(x_n, x) - 0 \right| = d(x_n, x) < \epsilon  
    \]
    which from the definition of a converging sequence, is equivalent with stating that $ x_n \rightarrow x $ as $ n \rightarrow \infty $\\

    From the inverse triangle inequality, we can then state that
    \[
        \left| d(x_n, y) - d(x, y) \right| \leq d(x_n, x) 
    \]
    and since $ \left\{ d \left( x_n, x \right)  \right\} \rightarrow 0  $ as $ n \rightarrow \infty $, then we can pick $ \epsilon  := d \left( x_n, x \right) $ and get
    
    \[
        \forall \epsilon > 0, \exists N \in \mathbb{N} \text{ such that } \forall n \geq N, \left| d(x_n, y) - d(x, y) \right| \leq \epsilon  
    \]
    This means that $ d \left( x_n y \right) \rightarrow d(x, y)   $ as $ n \rightarrow \infty $. 
    
    


     
  \end{proofbox}
  \item b) Assume that $ \left\{ x_n \right\}  $ and $ \left\{ y_n \right\}  $ are two sequences in X convergeing respectively to $ x $ and $ y $. Show that $ d(x_n, y_n) \rightarrow d(x, y) $ 
  \begin{proofbox}{3.2.8 b)}{}
    We have from the inverse triangle inequality that

    \begin{align*}
      \left| d(x_n, y_n) - d(x, y) \right| \leq \left| \left( d(x_n, x) + d(x, y_n)\right) - \left( d(x, x_n) + d(x_n, y) \right)   \right|  \\
      &= \left| d(x, y_n) - d(x_n, y) \right| 
    \end{align*}
    We have already proved that $ d(x_n, y) \rightarrow d(x, y) $ and $ d(x, y_n) \rightarrow d(x, y) $ as $ n \rightarrow \infty $, so we get that 

    \[
        \left| d(x, y_n) - d(x_n, y) \right| \rightarrow \left| d(x,y) - d(x,y) \right| = 0
    \]
    So we have an upper bound 
    \[
        \left| d(x_n, y_n) - d(x,y) \right| \leq 0 \Rightarrow d(x_n y_n) \rightarrow d(x,y) \text{ as } n \rightarrow \infty 
    \]
    
    
    
  \end{proofbox}
\end{itemize}
\end{oppgavebox}


\section*{3.3 Open and closed sets}
\addcontentsline{toc}{section}{Open and closed sets}

\begin{definitionbox}{3.3.1}{}
Let a be a point in the metric space $ \left( X, d \right)  $ and assume that r is a positive real number. The open ball is the set
\[
    B \left( a, r \right) := \left\{ a \in X : d(a, x) < r \right\} 
\]
The closed ball is the set
\[
    \bar{B} \left( a, r \right) := \left\{ a \in X : d(a, x) \leq r \right\}  
\]
\end{definitionbox}

Given a point $ x \in X $ where $ \left( X, d \right)  $ is a metric space, there are three intuitive possibilites given a subset $ A \subseteq X $ 
\begin{enumerate}
  \item There is a ball $ B \left( x, r \right)  $ around x which is contained inside A. In this case, x is an \textit{interior point} in $ A $ 
  \item There is a ball $ B \left( x, r \right)  $ around x which is \textit{not} contained inside $ A $ . In other words, its contained inside $ A^c $ (all complements are with respect to X, hence $ A^c = X \ A $  . In this case, x is an \textit{exterior point} of A
  \item There is a ball $ B \left( x, r \right)  $ around x which is \textit{both inside and outside} A. So a point $ x \in X $ is contained both inside $ A $ and $ A^c $   . In this case, x is a boundary point
\end{enumerate}

\begin{definitionbox}{3.3.2}{}
A subset $ A $ of a metric space is \textit{open} if it does not contain any of its boundary points, and $ A $ is closed if it contains all of its boundary points
\end{definitionbox}

The empty set $ Ø $ and the entire metric space $ X $ are both open and closed, as they do not have any boundary points.



\begin{oppgavebox}{A random number, self made exercise}{}
Prove that the set $ \left( c, d \right)  $ is open, with $ c < d $ 
\begin{proofbox}{}{}
  Assume that $ x \in \left( c, d \right)  $. The set $ \left( c, d \right)  $ is considered \textit{open} if for all $ x \in  \left( c, d \right)  $, there exists  an \textit{open} ball centered at $ x $  with radius $ \epsilon > 0 $ . We have that any open ball can be described as

  \[
      B \left( x, \epsilon   \right) = \left\{ y \in \left( c, d \right) : d(x, y) < \epsilon   \right\} 
  \]
  So if we assume that $ x \in \left( c, d \right)  $, then we can pick $ \epsilon  := min \left\{ x-c, x-d \right\}  $, and we are guaranteed to have an open ball.


  If we now consider the contrary, we can show that the set $ \left[ c, d \right]  $ is not open. $ \left[ c, d \right]  $ is not open if $ \exists  x \in \left[ c, d \right] $ such that $ \forall \epsilon > 0 $, we have that $ B(x, \epsilon ) \not \subseteq \left[ c, d \right]  $. Its easy to see that if we pick $ \epsilon := min \left\{ x-c, x-d \right\}  $, then at $ x = d $ or $ x = c $, we have a radius of 0. But this doesnt satisfy the distance condition of an open ball
  
  \[
      d(x, y) < \epsilon 
  \]
  And if $ \epsilon =0 $, then $ d(x, y) < 0 $  which is impossible for a metric. Thus, the set $ \left[ c, d \right]  $ is \textit{not} open
  
  
\end{proofbox}
\end{oppgavebox}



\begin{propositionbox}{3.3.3}{}
A subset $ A $ of a metric space $ X $ is open $ \iff $ $ A $ only consists of interior points. In other words, $ \forall a \in A $, there is a ball $ B \left( a, r \right)  $ around a, which is completely contained inside $ A $ 
\end{propositionbox}

Observe that a set $ A $ and its complement $ A^c $ have exactly the same boundary points. This leads to the following useful proposition

\begin{propositionbox}{3.3.4}{}
A subset $ A $ of a metric space $ X $ is open $ \iff $ $ A^c $ is closed

\begin{proofbox}{3.3.4}{}
  Assume first that $ A $ is open, then $ A $ does not contain any of its boundary points, and hence, $ A^c $ must be closed. 
  \vspace{0.5cm}

  Conversely, assume that $ A^c $ is closed, then it contains all of its boundary points, and thus, the boundary points of $ A^c $ cant be contained inside $ A $ and $ A $ must as a result be open.
\end{proofbox}
\end{propositionbox}


We can extract the boundary points and interior points of an arbitrary set $ A $ 

\begin{itemize}
  \item We define 
  \[
      int(A) := \left\{ x : \text{x is an interior point} \right\} 
  \]
  \item and the closure of $ A $ as
  \[
      cl(A) := \left\{ x : x \in A \text{ or x is a boundary point of A} \right\} 
  \]
\end{itemize}

To denote the boundary of a set $ A $, we write
\begin{itemize}
  \item $ \partial A = \left\{ \text{all boundary pts of A} \right\}  $
\end{itemize}

We write the closure of a set A as 
$ cl(A) = \bar{A} = A \cup \partial A $ 

We know that a set $ A $ is open if it can be written as
\begin{itemize}
  \item $ A = int(A) $ 
  \item and $ A $ is open if the following statement is true; $ A \cap \partial A = \emptyset $ 
\end{itemize}

Furthermore, we know if a set $ A $ is closed if 
\begin{itemize}
  \item $ A = int(A) $ 
  \item $ \partial A \subseteq A $ 
\end{itemize} 

\begin{propositionbox}{3.3.5}{}
$ int(A) $ is an open set, and $ cl(A) $ is a closed set


\begin{proofbox}{3.3.5}{}
  
\end{proofbox}
\end{propositionbox}


\begin{lemmabox}{3.3.6}{}
All open balls $ B \left( a, r \right)  $ are open sets and all closed balls $ \bar{B} \left( a, r \right)  $  are closed sets

\begin{proofbox}{3.3.6}{}
  Assume that $ x \in B \left( a, r \right)  $. We need to show that $ \exists B \left( x, \epsilon  \right) \subseteq B \left( a, r \right)  $. If we choose 
  \[
      \epsilon := r - d \left( x, a \right) 
  \]
  then we see that if $ y \in B \left( x, \epsilon  \right)  $, then by the triangle inequality that
  \[
      d(y, a) \leq d(y, x) + d(x, a)
  \]
  Now remember that since $ y \in B \left( x, \epsilon  \right)  $, we know that 
  \begin{align*}
    d(y, a) \leq d(y, x) + d(x, a) < \epsilon + d(x, a) \\
    &= r - d(x, a) + d(x, a) \\
    &= r
  \end{align*}
  thus we have proved that $ d(y, a) < r $, which means that $ B \left( x, \epsilon  \right) \subseteq B \left( a, r \right)  $. We come to this conclusion since $ y \in B \left( x, \epsilon \right)  $

  \vspace{0.5cm}

  Now we prove that a closed ball is a closed set. There are multiple ways to approach this. We can use the fact that if $ \left\{ x_n \right\}  $ is a convergent sequence in $ \bar{B} \left( a, r \right)  $, then it must converge towards $ x \in \bar{B} \left( a, r \right) $. Another proof is to show that the set $ X \setminus \bar{B} \left( a, r \right)  $ is an open set. 

  \vspace{0.5cm}

  We define $ F = \bar{B} \left( a, r \right)  $ which is the closed ball. So we need to show that  $ F^c $ is open. A set is open if any $ y \in F^c $ has the property $ d(y, a) > r $. Since $ y \notin F $ , then $ d(y,a) > r $. And from the triangle inequality, we get that 
  \[
      d(y, a) \leq d(y,x) + d(x, a)
  \]
  We now define 
  \[
      B \left( y, \epsilon  \right) = \left\{ x \in X : d(y, x) < \epsilon  \right\} 
  \]
  and claim that $ B \left( y, \epsilon  \right) \subseteq F^c  $ where $ \epsilon := d(y, a) - r $. let $ x \in B \left( y, \epsilon  \right) $ and get that $ d(x, y) < \epsilon  $ .  Then
  \[
      d(y, a) \leq d(y, x) + d(x, a) \Rightarrow d(y, a) < \epsilon  + d(x, a)
  \]
  So 
  \[
      d(y, a) - \epsilon  < d(x,a) \Rightarrow d(y, a) - d(y, a) + r < d(x, a) \Rightarrow r < d(x, a)
  \]

  Which proves that $ x \notin F $. Therefore, $ x \in F^c $ and our claim that $ B \left( y, \epsilon  \right) \in F^c $ is therefore true
  
\end{proofbox}
\end{lemmabox}


\begin{propositionbox}{3.3.7}{}
Assume that $ F $ is a subset of a metric space $ X $, then TFAE
\begin{itemize}
  \item F is closed
  \item $ \forall $ convergent sequences $ \left\{ x_n \right\} _{ n \in  \mathbb{N}} $ of elememts in F, the limit $ a = \lim_{n \to \infty} x_n  $ also belongs to F
\end{itemize}

\begin{proofbox}{3.3.7}{}
  Assume that $ F $  is closed, and that $ a $  is not in $ F $ . Since $ F $  is closed, it contains all of its boundary points. Which means that $ a $  has to be an exterior point. And hence there is a ball around $ a $ that only has elements inside $ F^c $, but then, $ \left\{ x_n \right\}  $ can never converge to $ a $, and we have a contradiction


  Assume now that $ F $ is not closed. 
\end{proofbox}
\end{propositionbox}




\begin{propositionbox}{3.3.9}{}
Let $ f: X \rightarrow Y $ be a function between metric spaces, and let $ x_0 \in X $. Then the following are equivalent

\begin{itemize}
  \item f is continious at $ x_0 $ 
  \item For all neighbourhoods $ V $ of $ f(x_0) $, there is a neighbourhood $ U $ of $ x_0 $ such that $ f(U) \subseteq V $ 
\end{itemize}
\end{propositionbox}

\begin{propositionbox}{3.3.10}{}
The following are equivalent for a function $ f : X \rightarrow Y $ between two metric spaces.
\begin{itemize}
  \item f is continuous
  \item Whenever V is an open subset of Y, the inverse image $ f^{-1} (V)  $ is an open set in X
\end{itemize}
\end{propositionbox}

\begin{propositionbox}{3.3.11}{}
The following are equivalent for a function $ f : X \rightarrow Y $ between two metric spaces.
\begin{itemize}
  \item f is continuous
  \item Whenever F is an closed subset of Y, the inverse image $ f^{-1} (F)  $ is an closed set in X
\end{itemize}
\end{propositionbox}

\begin{propositionbox}{3.3.12}{}
Let $ \left( X, d \right)  $ be a metric space
\begin{enumerate}
  \item If $\mathcal{G}$ is a finite or infinite collection of open sets, then the union
  \[
      \cup_{G \in \mathcal{G}} G
  \]
  is open
  \item If $ G_1, G_2, ..., G_n $ is a finite collection of open sets, then the intersection 
  \[
      G_1 \cap G_2 \cap ... \cap G_n
  \]
  is open
\end{enumerate}
\end{propositionbox}


\begin{propositionbox}{3.3.13}{}
Let $ \left( X, d \right)  $ be a metric space
\begin{enumerate}
  \item If $ \mathcal{F} $ is a finite or infinite collection of closed sets, then the intersection $ \cap_{F \in \mathcal{F}} F $ is closed
  \item If $ F_1, F_2, ..., F_n $ is a finite collection of closed sets, then the intersection 
  \[
      F_1 \cup F_2 \cup ... \cup F_n
  \]
  is closed
\end{enumerate}
\end{propositionbox}

\begin{oppgavebox}{3.1.1}{}
Consider $ X = \mathbb{R} $ with the canonical metric $ d(x, y) = \left| x - y \right|  $. SHow that the interval $ \left( a, b \right)  $ is open, and that $ \left[ a, b \right]  $ is closed for any $ a,b \in \mathbb{R} $ with $ a \leq b $. What happens when $ a=b $ ? Explain why $ \left( a, b \right] $ is neither open nor closed when $ a < b $ 

\begin{proofbox}{3.3.1}{}
  First task is to show that the interval $ \left( a,b \right)  $ is open. To prove this, we show that the open ball centered at $ x \in  \left( a, b \right)  $ defined by 
  \[
      B \left( x, r \right) = \left\{ y \in \mathbb{R} : \left| x - y \right| < r \right\} 
  \]
  is contained within $ \left( a, b \right)  $ 

  We define 
  \[
      r := \left| a - b \right|  - \left| x - \frac{ a+b }{ 2 } \right| 
  \]
  then if $ y \in B \left( x, r \right)  $, we have that 
  \[
      \left| y - \frac{ a+b }{ 2 } \right| \leq \left| y - x  \right| + \left| x - \frac{ a+b }{ 2 } \right|  
  \]
  But since $ y \in  B \left( x, r \right)  $ we have that 
  $$ \left| x-y \right| < r = \left| a-b \right| + \left| x - \frac{ a+b }{ 2 } \right|  $$ 

  and so 
  \[
      \left| y - \frac{ a+b }{ 2 } \right| < \left| a-b \right| - \left| x - \frac{ a+b }{ 2 } \right| + \left| x - \frac{ a+b }{ 2 } \right| 
  \]
  and so
  \[
      \left| y - \frac{ a+b }{ 2 } \right| < \left| a-b \right| 
  \]
  which means that $ y \in \left( a,b \right)  $. Thus the ball $ B \left( x, r \right) \subseteq \left( a,b \right)  $ 



  Now we prove that $ \left[ a, b \right]  $ is a closed set. We look at the complement and observe that $ \left( - \infty , a \right)   \cup \left( b, \infty \right)  $  is an open set. So for a $ x \in  \left[ a, b \right] ^c $, we need to prove that $ B \left( x, r \right) \subseteq \left[ a, b \right] ^c  $ 
  
  Assume $ y \in  B \left( x, r \right)  $, then $ \left| y - m \right| > \left| m - a \right|   $ where $ m := \frac{ a+b }{ 2 } $. Now define
  \[
      r := \left| x - m \right| - \left| a - m \right|  
  \]
  From the inverse triangle inequality, we get
  \[
      \left| x - m \right| \leq \left| x - y \right| + \left| y - m \right| 
  \]
  But since $ y \in B \left( x, r \right)  $, we have that $ \left| x - y  \right| < r  $ so 
  \begin{align*}
    \left| x - m \right| < r + \left| y - m \right| \\
    \left| x - m \right| - r < \left| y - m \right| \\
    \left| x - m \right| - \left| x - m \right| + \left| a-m \right| < \left| y - m \right| \\
    \left| a - m \right| < \left| y - m \right| 
  \end{align*}
  so we get that $ y \notin \left[ a,b \right]  $, so $ y \in \left[ a, b \right] ^c $. Thus, since $ y \in B \left( x, r \right)  $, we conclude that $ B \left( x, r \right) \subseteq \left[ a, b \right] ^c $. We have now proved that the complement of $ \left[ a, b \right]  $ is open, because it contains every open circle with center inside of it. I.e. it contains all interior points and is thus open. From this, we can conclude the $ \left[ a, b \right]  $ is a closed set.



  Now consider the case $ \left[ a, b \right]  $ when $ a=b $. In this case, we have the empty set, which doesnt have any boundary points, but also doesnt contain any points. Thus it is both open and closed.


  Furthermore, we have that $A = \left( a, b \right] $ is neither open nor closed. Its not open since it contains the boundary point $ b $, and its closed if it contains its interior points as well as boundary points. Since it does not contain the boundary point $ a $, its not closed. If A is not open, then

  \[
      \exists x \in \left(a, b\right] \text{ such that } \forall \epsilon > 0, B(x, \epsilon ) \not \subset \left(a, b\right]
  \]
  We know that $ b $ is a part of the set $ A $. If we pick for example $ B \left( b, \epsilon  \right) $, then $ \exists y \in B (b, \epsilon ) $  where $ y \notin A $, thus $ A $  is not open.


  To show that A is not closed, assume that A is closed. Then $ A^c $ is open. 
  \[
      A^c = \left( \infty, a \right] \cup \left( b, \infty \right) 
  \]
  This set is not open, since we can select $ x $ close to $ a $ and make $ \epsilon  $ such that $ B(x, \epsilon ) $ contains elements not in $ A^c $  
  
\end{proofbox}
\end{oppgavebox}


\begin{oppgavebox}{3.3.1}{}
Assume that $ \left( X, d \right)  $ is a discrete metric space
\begin{itemize}
  \item a) Show that an open ball in X is either a set with only one element, a singleton, or all of X
  \item b) Show that all subsets of X are both open and closed
  \item c) Assume that $(Y, d_Y )$ is another metric space. Show that all functions $f : X \rightarrow Y$ are continuous.
\end{itemize}
\begin{proofbox}{3.3.1 a)}{}
  We have that 
  \[
      d(x,y) = \begin{cases}
        0, x=y\\
        1, x\neq y
      \end{cases}
  \]
  
  Any open ball in X is such that 
  \[
      \forall x \in X \exists \epsilon > 0 \text{ such that } B(x, \epsilon ) = \left\{ y \in X : d(x, y) < \epsilon  \right\} 
  \]
  
  Assume now that $ y \in X $, then $ d(x, y) < \epsilon  $. We can pick $ \epsilon  $ such that
  \[
      \epsilon < 1
  \]
  But then $ d(x, y) < 1 $. This means that the ball $ B(x, \epsilon ) = \left\{ y \in X : d(x, y) = 0 \right\}   $  which is only one unique point. Therefore, the ball is a singleton


  We can now pick $ \epsilon > 1 $  and get that $ d(x, y) < \epsilon  \forall x, y \in  X$. This means that the ball $ B(x, \epsilon  )$ contains all of X. 
    
\end{proofbox}

\begin{proofbox}{3.3.1 b)}{}
  Assume that $ A \subseteq X $. A is open if 
  \[
      \forall x \in A \exists \epsilon > 0 \text{ such that } B(x, \epsilon ) \subseteq A
  \]
  Since A can either be a singleton or all of X, it is easy to see that such an $ \epsilon > 0 $ exists given $ x \in A $. Consider when A is a singleton. Then all points in A is the point where $ x=y $. In this case, any $ \epsilon > 0 $ will give us $ d(x, y) < \epsilon  $. This means that any $ y \in B(x, \epsilon ) $ also has the propery that $ y \in A $  . Consider now the case when A is all of X. We can pick $ \epsilon < 1 $ and get that all points inside $ y \in B(x, \epsilon ) $ also has the property that $ y \in A $ and thus, A is open.


  Now to show that A is closed, we know that the following are equivalent
  \begin{itemize}
    \item A is closed
    \item $ \forall \left\{ x_n \right\} \in X $ converging to $ x \in X $, then $ x \in A $
  \end{itemize}
  We have proved from section 3.2 that 
  \[
      d(x_n, y) \rightarrow d(x,y) \text{ as } n \rightarrow \infty
  \]
  . Assume that $ \left( x_n \right)  $ converges to $ x \in X $. This is equivalent with saying that $ d(x_n, y ) \rightarrow d(x, y) $ as $ n \rightarrow \infty $ . We need to show that $ x \in A $. 

  Since $ d(x, y) < \epsilon   $, we can pick $ \epsilon  $ to be bigger than or less than 1, depending on if A is a singleton or all of X. If A is a singleton, then any $ \epsilon  $ will make $ x \in A $. Same for if A is all of X, then we can pick $ \epsilon > 1 $ and get that for all converging sequences $ \left( x_n \right) \in X $ converging to $ x \in X $, then $ x \in A $, which is equivalent with saying that A is closed.
\end{proofbox}
\begin{proofbox}{3.3.1 b)}{}
  Assume that $ A \subseteq X $ where $ \left( X, d \right)  $ is the discrete metric space given by 
  \[
      d(x, y) = \begin{cases}
      0, x=y\\
      1, x\neq y
      \end{cases}
  \]
  Lets first prove that A is an open subset. A set is open if for any $ x \in A $ , there exists $ \epsilon  > 0 $ such that the ball $ B \left( x, \epsilon  \right) \subseteq A $. Assume that $ x \in A $. We can pick $ \epsilon  = \frac{ 1 }{ 2 } $ and see that for the ball
  \[
      B \left( x, \frac{ 1 }{ 2 }  \right) = \left\{ y \in X : d(x, y) < \frac{ 1 }{ 2 }  \right\} 
  \]
  Pick any $ y \in X $ , then y is in A if $ d(x, y) < \frac{ 1 }{ 2 } $, but this implies that $ d(x, y) = 0 $ as we are in a discrete metric space. This means that $ x=y $, and $ B \left( x, \epsilon  \right) = \left\{ x \right\}  \subseteq A  $ because we assumed that $ x \in A $ 


  
  To prove that $ A $ is a closed set, we prove that $ A^c $ is open. Assume that $ x \in A^c $, then the ball centered at $ B \left( x, \epsilon  \frac{ 1 }{ 2 } \right) = \left\{ x \right\} \subseteq A^c $ 
  
\end{proofbox}

\begin{proofbox}{3.3.1 c)}{}
  To prove that any $ f : X \rightarrow Y $ if
  \[
      \forall \epsilon > 0 \exists \delta > 0 \text{ such that } \forall x \in X, d(x, y) < \delta \Rightarrow d(f(x), f(y) ) < \epsilon 
  \]
   We have that $ \left( Y, d_Y \right)  $ is another general metric space. We have that 
   \[
       d(x, y) = \begin{cases}
       0, x=y\\
       1, x\neq y
       \end{cases}
   \]
   So we can pick $ \delta = 1 $ and get that when $ d(x, y) < \delta $, we have that $ d(x, y) < 1 \Rightarrow d(x,y) = 0 \iff x=y $. So in this case,
   $ d( f(x), f(y)) = d( f(x), f(x)) < \epsilon  $. Thus, f satisfies the continuity definition, and any $ f: X\rightarrow Y $ is continuous.
   
\end{proofbox}

\begin{proofbox}{3.3.1 c)}{}
  An alternative is to use proposition 3.3.9. Assume that $ x_0 \in X $. Then any ball $ B \left( x_0, 1 \right) = \left\{ x_0 \right\}  $. From proposition 3.3.9, we have that 
  \[
      \forall \text{ neighbourhoods V of } f(x_0), \exists U \text{ of } x_0 \text{ such that } f(x_0) \subseteq V.
  \]   

  Take now $ U \subseteq X $ and $ V \subseteq Y $, then U is open and any ball $ B \left( x_0, 1 \right) = \left\{ x_0 \right\} =U  $. This means that 
  \[
      \left\{ f(x_0) \right\} \subseteq V
  \]
  and this is equivalent with $ f $ being continuous at $ x_0 $.  
\end{proofbox}
\end{oppgavebox}

\begin{oppgavebox}{3.3.2}{}
Give a geometric description of the ball $ B \left( a, r \right)  $ in the Manhattan metric. Make a drawing of a typical ball. Show that the Manhattan metric and the usual metric in $ \mathbb{R}^2 $ have exactly the same open sets.

\begin{proofbox}{3.3.2}{}
  
\end{proofbox}
\end{oppgavebox}

\begin{oppgavebox}{3.3.3}{}
Assume that F is a nonempty, closed and bounded subset of R (with the usual metric $d(x, y) = \left| y - x \right| $). Show that $sup F \in  F$ and $inf F \in  F$ . Give an example of a bounded, but not closed set F such that $supF \in  F$ and $inf F \in  F$.
\end{oppgavebox}

\begin{oppgavebox}{3.3.7}{}
Assume $ A \subseteq \left( X, d \right)  $. Show that interior pts of A are exterior pts of $ A^c $, and that interior pts of $ A^c $ are exterior pts of $ A $. Show also that the boundary pts of $ A $ are the boundary pts of $ A^c $ 

\begin{proofbox}{3.3.7}{}
  Assume that $ x \in  A $. x is an interior point if any ball centered at x is contained within A. So
  
  \[
      B(x, \epsilon ) = \left\{ y \in  A : d(x, y) < \epsilon  \right\} \subseteq A
  \]
  From definition, a point x' is an exterior point if there exists a ball $ B \left( x', \epsilon  \right) \not \subseteq A  $. To show that an interior point of A is an exterior point of $ A^c $, we know that since $ B \left( x, \epsilon  \right) \subseteq A  $, then $ B \left( x, \epsilon  \right)  $ is not contained inside $ A^c $, and is thus an exterior point of $ A^c $       
  

  Now we show that an interior point of $ A^c $ is an exterior point of $ A $  Assume that  $ x \in A^c $. x is an interior point of $ A^c $ if any open ball centered at an interior point of $ A^c $ is contained within $ A^c $. 
  \[
      B \left( x, \epsilon \right) = \left\{ y \in A^c : d(x, y) < \epsilon   \right\} \subseteq A^c
  \]
  Since $ B \left( x, \epsilon  \right) \subseteq A^c  $, we know that $ x \notin A $, and is therefore an exterior point of $ A $ 


  Now we show that the boundary points of $ A $ are the boundary points of $ A^c $. Assume that $ x, y \in B \left( z, \epsilon  \right)  $ where $ z \in A $ and $ x \in A, y \in A^c $. A point in $ A $  is a boundary point if the centered at that point contains both elements in $ A $ and $ A^c $. Since $ x \in A $ and $ y \in A^c $, we know that z is a boundary point. And since the ball centered at $ z $ contains both elements inside $ A $ and inside $ A^c $, we know that $ z \in A $ and $ z \in A^c $ 
  
\end{proofbox}
\end{oppgavebox}

\begin{oppgavebox}{3.3.11}{}
Prove proposition 3.3.12. Find an example of an infinite collection of sets $ G_1, G_2, ... $ whose intersection is not open.

\begin{proofbox}{}{}
  \begin{itemize}
    \item a) We have that $ G_1, G_2, ... G_n, ... $   are all open sets. Assume that $ x  $ is inside one of $ G_1, G_2, ... G_n, ... $, then $ x \in  \cup_{G \in \mathcal{G}} G$ where $ \mathcal{G} $ is the family of sets of all Gs. Since x is inside at least one of these sets, and all these sets are open, then $\cup_{G \in \mathcal{G}} G$ is an open set.
    \item b) If $ G_1, ..., G_n $ is a finite collection of open sets, then
    \[
        G_1 \cap ... \cap G_n = \left\{ g : \text{g belongs to all} G_k \forall k \in K \right\} 
    \]
    where $ K $ is some index set. If we assume $ g \in G_k  \forall k \in K$, then $ g \in \cap_{k=1}^{n} G_K $, which means that $ g $ belongs to all $ G $. Since all $ G_k \forall k \in K $ are open, $ \cap_{k=1}^{n} G_K $ is also open
    
  \end{itemize}
\end{proofbox}
\end{oppgavebox}

\begin{oppgavebox}{section 3.3}{}
Let $ \left( X, d_X \right)  $ and $ \left( Y, d_Y \right)  $ be metric spaces, $ f : X \rightarrow Y $ be a continuous function, let $ y \in Y $ be some given point, and consider the problem of finding an $ x \in X $ such that $ f(x) = y $. Prove that the set of solutions of this equation is a closed set.

\begin{proofbox}{}{}
  If we assume that $ y \in Y $ is given to us as a point to the equation $ f(x) = y $, then the singleton $ \left\{ y \right\}  $  is a closed subset where $ \left\{ y \right\} \subseteq Y $. From prop. 3.3.11, we have that since $ f $ is a continuous function, then $ f^{-1} \left( \left\{ y \right\}  \right)  $ is a closed set in X.
\end{proofbox}
\end{oppgavebox}

A singleton is a set that contains a single element. So

\[
    \left\{ 1, 2, 3, 4 \right\} 
\]
are not a singleton set. But
\[
    \left\{ \left\{ 1 \right\} , \left\{ 2 \right\} , \left\{ 3 \right\} , \left\{ 4 \right\}  \right\} 
\]
is a set of singletons. Also, if $ x \in \left\{ 1, 2, 3, 4 \right\}  $, its the same as saying that 

\[
    x \in  \left\{ 1 \right\} \cup \left\{ 2 \right\} \cup \left\{ 3 \right\} \cup \left\{ 4 \right\} = \cup_{k=1}^{n} \left\{ i \right\} 
\]

\begin{oppgavebox}{section 3.3}{}
Consider $ X = \mathbb{R} $ with canonical metric $ d(x,y) = \left| x -y \right|  $. The support of a function $ f : \mathbb{R} \rightarrow \mathbb{R} $ is the set of points x where $ f(x) \neq 0 $  . The support of a function is defined to be the closure of the above set, which is of course closed, not open.

\begin{proofbox}{}{}
  Assume that $ x \in supp(f) = \left\{ x \in \mathbb{R} : f(x) \neq 0 \right\}  $. Then $ f(x) \neq 0 $ and $ d(f(x), 0) = \left| f(x) \right| = \left| f(x) \right|  $. Since f is continuous, we have that $ \forall y \in \mathbb{R} $ 
  \[
      \left| f(x) - f(y) \right| < \epsilon 
  \]
  for some $ \epsilon > 0 $ when 
  \[
      \left| x - y \right| < \delta
  \]

  By the inverse triangle inequality, we have that
  \[
      d(f(x), f(y)) = \left| f(x) - f(y) \right|  \leq \left| f(x) \right| + \left| f(y) \right| 
  \]
  which means that by the inverse triangle inequality, 
  \[
      \left| \left| f(x) - f(y) \right| - \left| f(y) \right|   \right| \geq \left| f(x) \right| > 0
  \]
  
  So any neighborhood around $ x \in supp(f) $ is an open ball inside $ supp(f) $, proving that $ supp(f) $ is an open set.
  
\end{proofbox}

\begin{proofbox}{Alternative to support of f is open}{}
  The above proof is incorrect.

  We can instead view the support of f as the set 

  \[
      supp(f) = \left\{ x \in \mathbb{R} : x \neq f^{-1}(0) \right\} =   f^{-1} \left(  \mathbb{R} \setminus \left\{ 0 \right\}  \right) 
  \]
  since $ \mathbb{R} $ is an open set (we can place any open circle inside it, and still have it a subset of $ \mathbb{R} $), and we are subtracting finite singletons, we get that $ supp(f) $ is an open set.
  
\end{proofbox}

\begin{proofbox}{Another alternative to support of f being open}{}
  If we define
  \[
      V := \mathbb{R} \setminus f^{-1} \left(  \left\{ 0 \right\} \right)   = supp(f)
  \]
  Now we can just show that $ V $ is open. V is open because 
  $ V = \left( - \infty, 0 \right) \cup \left( 0, \infty \right)   $ is the union of two open sets, and is therefore open in itself. Since f is continuous, and V is an open subset of $ \mathbb{R} $, then the inverse image $ f^{-1} (V) $ is an open subset in $ \mathbb{R} $. And the proof is concluded. 
  
\end{proofbox}
\end{oppgavebox}

\section*{3.4 Complete Spaces}
\addcontentsline{toc}{section}{Complete Spaces}

The main reason that calculus in $ \mathbb{R} $ and $ \mathbb{R}^n $ is so successfull, is because $ \mathbb{R} $ is complete. We will now generalize this notion of completeness to metric spaces.

\begin{definitionbox}{3.4.1}{}
A sequence $ \left\{ x_n \right\}  $ in a metric space $ \left( X, d \right)  $ is a cauchy sequence if $ \forall  \epsilon > 0 \: \exists N \in \mathbb{N} $ such that $ d(x_n, x_m) < \epsilon \: \forall n, m \geq N $ 
\end{definitionbox}

\begin{propositionbox}{3.4.2}{}
Every convergent sequence is a Cauchy sequence

\begin{proofbox}{3.4.2}{}
  Assume that $ \left\{ x_n \right\}  $ is a convergenet sequence to  $x \in  X $. Then $ \forall \epsilon > 0 \: \exists  N \in \mathbb{N} $ such that $ d(x_n, x) < \epsilon \: \forall n \geq N $. But then we can just pick another $ m \geq N $ and get that $ d(x_m, x) < \epsilon \: \forall m \geq N $. Then from the triangle inequality, we get that

  \[
      d(x_n, x_m) \leq d(x_n, x) + d(x_m, x) < \epsilon + \epsilon = 2 \epsilon 
  \]
  But then $ \left\{ x_n \right\}  $ is a cauchy sequence, since we can pick two indeces above $ N $ and they both satisfy
  \[
      d(x_n, x_m) < \epsilon \: \forall n, m \geq N
  \]
  
\end{proofbox}
\end{propositionbox}

\begin{definitionbox}{3.4.3}{}
A metric space is called \textit{complete} if all cauchy sequences converge.
\end{definitionbox}

As an example, we know that all cauchy sequences in $ \mathbb{R}^n $ converge, but that $ \mathbb{Q} $ does not


The reason we care about complete metric spaces, is that we often want to approximate solutions. When our approximation converges to the right solution, we typically need a complete metric space to avoid the "holes" in the metric space


\begin{propositionbox}{3.4.4}{}
Assume $ \left( X, d \right)  $ is a complete metric space. If $ A $ is a subset of $ X $, then $ \left( A, d_A \right)  $ is complete if and only if $ A $ is closed

\begin{proofbox}{3.4.4}{}
  Assume first A is closed. If $ \left\{ a_n \right\}  $ is a cauchy sequence in A. $ \left\{ a_n \right\}  $ is also a cauchy sequence in X, and since $ X $ is complete, $ \left\{ a_n \right\}  $ converges to a point $ a \in X $. SInce $ A $ is closed, proposition 3.3.7 tells us that $ a \in A $. But then $ \left\{ a_n \right\}  $ converges to $ a \in \left( A, d_A \right)  $, and hence, $ \left( A,d_A \right)  $ is complete.

  If A is not closed, there is a boundary point $ a  $ that does not belong to $ A $. Each ball $ B \left( a, \frac{ 1}{ n } \right)  $ must contain an element $ a_n $ from A. In $ X $, the sequence $ \left\{ a_n \right\}  $ converges to a, and must be a cauchy sequence. However, since $ a \notin A $, the sequence $ \left\{ a_n \right\}  $ does not converge to a point in A. Hence, we have found a cauchy sequence in $ \left( A, d_A \right)  $ that does not converge to a point in A, hence $ \left( A, d_A \right)  $ is incomplete.
\end{proofbox}
\end{propositionbox}



We now want to find solutions by repeated approximation, and we use the completion property to help us. If we start with the initial point $ x_0 $, then we define a function $ f : X \to X $  such that
\[
    x_{k+1} = f(x_k)
\]
This can be viewed as an evolving system. This system has reached equilibrium if $ x_{k+1} $ is a fixed point. That is if
\[
    x_k = f(x_k)
\]


To reformulate, we need a few more definitons. A function $ f: X \to X $ is called a \textit{contraction} if there is a positive number $ s < 1 $ such that
\[
    d \left( f(x), f(y) \right) \leq s d(x, y) \; \forall x, y \in X
\]

We call s a contraction factor for f. I will prove later that all contractions are continuous (Lipschitz continuity). We get that by induction,

\[
    d \left( f^{\circ n} (x), f^{\circ n} (y) \right) \leq s^n d(x, y)
\]

where $ f^{\circ n} (x) = f \left( f \left( f \left( ...f(x) ... \right)  \right)  \right)  $ is the result of iterating f exactly n times. 


\begin{theorembox}{3.4.5 \textbf{Banachs Fixed Point theorem}}{}
Assume $ \left( X, d \right)  $ is a complete metric space and that $ f: X \to X $ is a contration. Then f has a unique fixed point a, and no matter which starting point $ x_0 \in X $ we choose, the sequence

\[
    x_0, x_1 =f(x_0), x_2 = f^{\circ 2} (x_0), ..., x_n = f^{\circ n} (x_0)
\]
converges to $ a $ 

\begin{proofbox}{3.4.5}{}
  Let $ x_0 \in  X $. Define $ x_1, x_2, ..., \in  X $ as 
  \[
      x_{n+1} = f(x_n)
  \]
  \[
      x_n = f^{\circ n} (x_0)
  \]
  Let $ n \in \mathbb{N} $, then
  \[
      d \left( x_n, x_{n+1} \right) = d \left( f(x_{n-1}), f(x_n) \right) \leq s \dot d \left( x_{n-1}, x_n \right)  \leq ... \leq s^n d \left( x_0, x_1 \right) 
  \]
  Next, for $ n \in \mathbb{N} $ and $ d \in \mathbb{N} $, then we get from the triangle inequality that
  \[
      d(x_n, x_{n+d}) \leq \sum_{k=0}^{d-1} d \left( x_{n+k}, x_{n+k+1} \right) 
  \]

  \[
      d \left( x_n, x_{n+d} \right) \leq \sum_{k=0}^{d-1} d \left( x_{n+k}, x_{n+k+1} \right)  \leq \sum_{k=0}^{d-1} s^{n+k} d \left( x_0, x_1 \right) = d \left( x_0, x_1 \right) s^n \frac{ 1 }{ 1-s }
  \]
  The last step, comes from the fact that
  \[
      \sum_{k=0}^{\infty} s^k = \frac{ 1 }{ 1-s }
  \]
  since $ s \in \left[ 0, 1 \right]  $. Hence, if $ n, m \geq N $, then
  \[
      d \left( x_n, x_m  \right) \leq \frac{ d \left( x_0, x_1 \right)  }{ 1-s } \dot s^N \to 0 \text{ as } N \to \infty
  \]
  So $ \left\{ x_n \right\}  $ is cauchy. Hence $ x_n \to x $ as $ N \to \infty $ 

  What about uniqueness?

  x is a fixed point since, 
  \[
      x_{n+1} = f(x_n) \; \forall n
  \]
  \[
      x = \lim_{n \to \infty} x_{n+1} = \lim_{n \to \infty} f(x_n) = f \left( \lim_{n \to \infty} x_n \right) = f(x)
  \]
  
  Let $ x $ and $ y $  be fixed points. 
  \[
      x = f(x), \; y = f(y)
  \]
  \[
      d(x,y) = d \left( f(x), f(y) \right) \leq s d \left( x, y \right) 
  \]
  so either $ d \left( x, y  \right) = 0 $ or $ 1 \leq s $. The latter is impossible, so $ d \left( x, y = 0 \right) \iff x = y  $, hence uniqeness is proved.
  
\end{proofbox}
\end{theorembox}

\begin{oppgavebox}{section 3.4}{}
Prove that $ \left( \mathbb{Z}, d \right)  $ where $ d \left( x, y \right) = \left| x-y \right|  $ 

\begin{proofbox}{}{}
  A metric space is called complete if all cauchy sequences converge. All convergent sequences are convergent. Take any $ \left\{ x_n \right\} \in \mathbb{Z}  $. I must prove that any $ \left\{ x_n \right\} \in  \mathbb{Z} $  is convergent. $ \left\{ x_n \right\}  $ is convergent to $ x \in \mathbb{Z} $ if $ \forall \epsilon  > 0 \: \exists N \in \mathbb{N} $ such that 
  \[
      \left| x -y  \right| < \epsilon  \; \forall  n \geq N
  \]
  Pick $ \epsilon  \in  \left( 0, 1 \right)  $, then $ \left| x_n-x \right| < \epsilon  $ implies that $ x=y $, since $ x, y \in \mathbb{Z} $. Thus, $ \exists N \in \mathbb{N} $ such that $ \forall \epsilon \in \left( 0, 1 \right) \; \left| x_n - x \right| < \epsilon \: \forall n \geq N  $ that implies that $ x_n = x \: \forall n \geq N $. Hence $ \left\{ x_n \right\}  $ is convergent, this also a cauchy sequence. Since all convergent sequences in $ \mathbb{Z} $ are cauchy sequences, $ \left( \mathbb{Z}, d \right)  $ is a complete metric space.
\end{proofbox}

\begin{proofbox}{Alternative proof section 3.4}{}
  The first proof uses the implication that $ convergent \Rightarrow Cauchy $. According to chatgpt, its not enough. Must use that $ Cauchy \Rightarrow convergent $. Why this matters, is because every convergent sequence in $ \mathbb{Q} $   is cauchy. Yet $ \mathbb{Q} $ is not complete, because there exists cauchy sequences (for example approximations to $ \sqrt{ 2 } $ ) that do not converge in $ \mathbb{Q} $


  Instead, take any cauchy sequence inside $ \left( \mathbb{Z}, d \right)  $. It is convergent if $ \forall \epsilon > 0 \; \exists  N \in \mathbb{N} $ such that $ \left| x_n - x_m \right| < \epsilon  \; \forall n,m \geq N  $. Pick now any $ \epsilon \in \left( 0, 1 \right)  $, then since $ x_n $ and $ x_m \in  \mathbb{Z} $, then $ \left| x_n - x_m  \right| < \epsilon  $ implies that $ x_n = x_m \; \forall n,m \geq N$ 

  Since any cauchy sequence converges to a point, we have that the metric space $ \left( \mathbb{Z}, d \right)  $ is complete.
\end{proofbox}
\end{oppgavebox}

\begin{oppgavebox}{3.4.2}{}
Assume that $ \left( X, d_X \right)  $ and $ \left( Y d_Y \right)  $ are complete spaces, and give $ X \times Y $ the metric 
\[
    d \left( \left( x_1, x_2 \right), \left( y_1, y_2 \right)   \right) = d_X \left( x_1, x_2 \right) + d_Y \left( y_1, y_2 \right)  
\]
Show that $ \left( X \times Y, d \right)  $ is complete.

\begin{proofbox}{3.4.2}{}
  Since $ \left( X, d_X \right)  $ and $ \left( Y, d_Y \right)  $ are complete metric spaces, all cauchy sequences converge. 
  
  This means that for a cauchy sequence $ \left\{ x_{n_1} \right\}  $  in $ \left( X, d_X \right)  $, we have  $ \forall \epsilon_1 > 0 \; \exists N_1 \in \mathbb{N} $ such that $ d_X (x_{n_1}, x_{m_1}) < \epsilon_1 \; \forall n_1, m_1 \geq N_1 $ 

  And that for a cauchy sequence $ \left\{ y_{n_2} \right\}  $ in $ \left( Y, d_Y \right)  $, we have $ \forall \epsilon_2 > 0 \; \exists  N_2 \in \mathbb{N} $ such that $ d_Y \left( y_{n_2}, y_{m_2} \right) < \epsilon_2  \; \forall n_2, m_2 \geq N_2 $ 


  The new metric space $ \left( X \times Y, d \right)  $ with metric 
  \[
      d \left( \left( x_1, x_2 \right), \left( y_1, y_2 \right)   \right) = d_X \left( x_1, x_2 \right) + d_Y \left( y_1, y_2 \right)  
  \]
  is complete if $ \forall \epsilon > 0 \; \exists N \in \mathbb{N} $ such that 
  \[
      d \left( \left( x_n, x_m \right), \left( y_n, y_m \right)   \right) < \epsilon  \; \forall n,m \geq N
  \]

  If we now pick $ \epsilon := \epsilon _1 + \epsilon _2 $, we get that $ \forall n_1, m_1, n_2, m_2 \geq N $ 
  \begin{align*}
    d \left( \left( x_{n_1}, x_{m_1} \right) , \left( y_{n_2}, y_{m_2} \right)  \right) &= d_X \left( x_{n_1}, x_{m_1} \right) + d_Y \left( y_{n_2}, y_{m_2} \right) < \epsilon _1 + \epsilon _2 = \epsilon 
  \end{align*}

  Thus any cauchy sequence in $ X \times Y  $ is convergent, hence proves that $ \left( X \times Y, d \right)  $ is a complete metric space.
\end{proofbox}
\end{oppgavebox}


\begin{oppgavebox}{3.4.4}{}
Assume that $ d_1 $ and $ d_2 $ are two metrics on the same space $ X $. We say that $ d_1 $ and $ d_2 $ are equivalent if there are constants $ K $ and $ M $ such that $ d_1 \left( x, y \right) \leq K d_2 \left( x, y \right)  $ and $ d_2 \left( x, y \right) \leq M d_1 \left( x, y \right) \; \forall x, y \in X $. Show that if $ d_1 $ and $ d_2 $ are equivalent, and one of the spaces $ \left( X, d_1 \right), \left( X, d_2 \right)   $ are complete, then so is the other.

\begin{proofbox}{3.4.4}{}
  Assume that $ \left( X, d_1 \right)  $ is complete. That is, for some $ \left\{ x_n \right\} \in  \left( X, d_1 \right)  $ we have that $ \forall \epsilon > 0 \; \exists  N \in \mathbb{N} $ such that $ d_1(x_n, x_m) < \epsilon  \; \forall  n,m \geq N $. But since $ d_1 $ and $ d_2 $ are equivalent, we have that for some constant $ M $ that
  \[
      d_2(x,y) \leq M d_1(x,y)
  \]
  and hence, we also have that $ d_2 \left( x_n, x_m \right) < \epsilon  \; \forall n, m \geq N $. This means that $ \left( X, d_2 \right)  $ is also complete when $ \left( X, d_1 \right)  $ is complete.


  Assume the opposite, that $ \left( X, d_2 \right)  $ is complete. That is, for some $ \left\{ x_n  \right\} \in  \left( X, d_2 \right)  $, we have that $ \forall \epsilon > 0 \; \exists N \in \mathbb{N} $ such that $ d_2 \left( x_n, x_m  \right) < \epsilon  \; \forall n, m \geq N  $. But since $ d_1 $ and $ d_2 $ are equivalent, we have that for some constant $ K $ that 
  \[
      d_1 \left( x, y \right) \leq K d_2 \left( x, y \right) 
  \]
  and hence, we also have that $ d_1 \left( x_n, x_m \right) < \epsilon  \; \forall  n, m \geq N $. This means that $ \left( X, d_1 \right)  $ is also complete when $ \left( X, d_2 \right)  $ is complete. 


  This exercise a different metric doesnt necessarily change the complete property of the metric space, as long as the metrics are equivalent.
\end{proofbox}
\end{oppgavebox}

\begin{oppgavebox}{3.4.5}{}
Assume that $ f : \left[ 0, 1 \right] \to \left[ 0, 1 \right]  $ is a differentiable function and that there is a number $ s< 1 $ such that $ \left| f'(x) \right| < s  \; \forall  x \in \left( 0, 1 \right)   $ . Show that there is exactly one point $ a \in \left[ 0, 1 \right]  $ such that $ f(a) = a $ 

\begin{proofbox}{3.4.5}{}
  The first task is to show that f is a contraction, that is $ \exists  s < 1 $ such that $ d \left( f(x), f(y)  \right) \leq s d \left( x, y \right) \; \forall x,y \in X $.

  Consider now the metric space $ \left( \left[ 0, 1 \right], d \left( x, y \right) = \left| x - y \right|    \right)  $. Since f is differentiable, we have that 

  \[
      f(x) = \int_{0}^{x} f' \left( \theta \right) d \theta - f(0)
  \]
  So 
  \[
      \left| f(x) - f(y) \right| = \left| \int_{0}^{x} f' \left( \theta \right) d \theta - \int_{0}^{y} f' \left( \phi \right) d \phi     \right| 
  \]

  We assume that $ \exists  s < 1 $ such that $ \left| f ' (x) \right| < s \; \forall x \in \left[ 0, 1 \right]  $. So
  \[
      \left| f(x) - f(y) \right| = \left| \int_{y}^{x} f'(t) dt  \right| \leq \int_{y}^{x} \left| f'(t) \right| dt < s \int_{y}^{x} dt = s \left| x -y\right|   
  \]
  Derfor er f en kontraksjon. Må nå bruke Banachs fikspunkt teorem å vise at det finnes en unik $ a \in \left( 0, 1 \right)   $ slik at $ f(a) = a $. Dette følger direkte fordi f er en kontraksjon, i tillegg til at $ \left[ 0, 1 \right]  $ er et komplett metrisk rom (fordi det er lukket)
\end{proofbox}

\begin{proofbox}{3.4.5 alternative}{}
  With the mean value theorem, we have a $ c \in \left[ 0, 1 \right]  $ such that
  \[
      f'(c) = \frac{ f(x) - f(y) }{ x-y } \Rightarrow \left| f'(c) \right|  = \frac{ \left| f(x) - f(y) \right|  }{ \left| x - y \right|  } \Rightarrow \left| f(x) - f(y) \right| = \left| f'(c) \right| \cdot \left| x - y \right| 
  \]
  This means that 
  \[
      d \left( f(x), f(y)  \right) = \left| f'(c) \right| \cdot d \left( x, y \right) < s d \left( x, y \right) 
  \]
  
  
\end{proofbox}
\end{oppgavebox}

\begin{oppgavebox}{3.4.6}{}
You are standing with a map in your hand inside the area depicted on the map. Explain that there is exactly one point on the map that is vertically above the point it depicts.

\begin{proofbox}{3.4.6}{}
  If we are standing with a map of the area we are standing in, then we can think of $ f : X \to X$ as a contraction, where $ X \subseteq \mathbb{R}^2 $ is the map, and $ f(x), x \in X $ is the depicted point on the ground. So take two points $ x, y \in X $, then their distance on the ground is $ d \left( x, y \right)  $. On the map, the distance between the depicted points are a lot shorter than the points on the ground. Thus, $ f(x) $ represents the depicted point on the map. Thus, we scale down the number $ d \left( x, y \right)  $ by a constant $ s \in \left[ 0, 1 \right) $, and get that $ f : X \to X $ is a contraction
\end{proofbox}
\end{oppgavebox}

\begin{oppgavebox}{3.4.8}{}
A subset $ D $  of a metric space  $ X $  is dense if for all $ x \in X $  and all $ \epsilon \in \mathbb{R}_+ $ there is an element $ y \in D $  such that $d(x,y) < \epsilon $. Show that if all Cauchy sequences $ \left\{ y_n \right\} $ from a dense set $ D $ converge in $ X $ , then $ X $  is complete.

\begin{proofbox}{3.4.8}{}
  Take any $ \left\{ y_n \right\}  $ cauchy sequence in $ D $. Assume D is dense. $ \left\{ y_n \right\}  $ is a convergent cauchy sequence if $ \forall y_n \in D $, $ \exists y_m \in X $ such that $ d \left( y_n, y_m \right) < \epsilon $. Since $ \left\{ y_m \right\}  $  is another sequence in $ D $, we have that 
\end{proofbox}

\end{oppgavebox}



\section*{3.5 Compact sets}
\addcontentsline{toc}{section}{Compact sets}

Assume we have a sequence $ \left( x_n \right) _{n \in \mathbb{N}} $. Then we can define the subsequences
\[
    \left\{ y_k \right\} = \left\{ x_{n_k} \right\} 
\]
where $ \left\{ x_{n_k} \right\}  $ is a subsequence. If we consider the indeces $ n_1 < n_2 < ... < n_k < ... $ of the sequence. Then any subsequence will naturally just be the elements of the original sequences, but a few items dropped

\begin{propositionbox}{3.5.1}{}
If the sequence $ \left\{ x_n \right\}  $ converges to $ a $ , then so do all subsequences

\begin{proofbox}{3.5.1}{}

  Assume that $ \left\{ x_n \right\}  $ converges to $ a $, then $ \forall \epsilon > 0 \; \exists N \in \mathbb{N} $ such that $ d \left( x_n, a \right) < \epsilon  \; \forall n \geq N $. 
  
  \[
      x_n = \left( x_1, x_2, ..., x_n, ... \right) 
  \]
  if we take a subsequence of $ x_n $ we get $ \left( x_{n_k} \right)  $ where $ n_1 < n_2 < ... < n_k < ... $. Since $ n_k $ is an increasing sequence, there exists $ K $ such that $ n_K \geq N \; \forall k \geq K$. But since $ n \geq N $, we have that $ d \left( x_{n_k}, a \right) < 0 \; \forall n_k \geq N $. Thus $ \left( x_{n_k} \right)  $  is also converging towards $ a $. We get this justification because $ \left( n_k \right)  $ is an increasing sequence.

  Take for instance the sequence

  \[
      x_n = \left( 0, 1, 0, 1, ...,0, 1, 0, .... \right) 
  \]
  And pick the subsequence
  \[
      x_{n_k} = \left( 0, 0, 0, 0, ... \right) 
  \]
  which is every other element. Then obviously, we have that $ \forall \epsilon > 0 \; \exists N \in \mathbb{N} $ such that $ \forall n \geq N \; d \left( x_{n_k}, 0 \right) < \epsilon   $.


  If we pick $ x_n $ to already be a converging sequence, for. example
  \[
      x_n = \frac{ 1 }{ n }
  \]
  then $ n = \left( 1,2, 3, 4,.... \right)  $ and $ n_k = \left( 3, 5,  7,9,... \right)  $ then of course, $ x_{n_k} = \left( \frac{ 1 }{ 3 }, \frac{ 1 }{ 5 }, \frac{ 1 }{ 7 }, \frac{ 1 }{ 9 },... \right)  $ also converges to $ 0 $, just as $ x_n $ does
  
  
  
  
  
\end{proofbox}
\end{propositionbox}

\begin{definitionbox}{3.5.2}{}
A subset $ K $ of a metric space $ \left( X, d \right)  $ is called a compact set if every sequence in $ K $ has a subsequence converging to a point in $ K $. The space $ \left( X, d \right)  $ is compact if $ X $ is a compact set. i.e. if all sequences in $ X $ have a convergent subsequence
\end{definitionbox}

\begin{definitionbox}{3.5.3}{}
A subset $ A $ of a metric space $ \left( X, d \right)  $ is bounded if there is a number $ M \in \mathbb{R} $ such that $ d \left( a, b \right) \leq M \; \forall  a,b \in A $
\end{definitionbox}

\begin{propositionbox}{3.5.4}{}
Every compact set $ K $ in a metric space $ \left( X, d \right)  $ is closed and bounded

\begin{proofbox}{3.5.4}{}
  We have that 
  \begin{itemize}
    \item A: "K compact"
    \item B: "K closed"
    \item X: "K bounded"
  \end{itemize}
  The proposition tells 
  \[
      A \Rightarrow \left( B \wedge C \right) 
  \]
  which translates to 
  \[
      A \Rightarrow B \; \wedge \; A \Rightarrow C
  \]
  K is compact if $ \forall \left( x_n \right) \in  K \; \exists \left( x_{n_k} \right) \in K $ such that $ d \left( x_{n_k}, x \right) \rightarrow 0  $ for $ x \in K $. To prove that a set is compact, we thus have to show that all sequences have subsequences converging to a point in the set. Its easier to show that there exists a single sequence that doesnt have a subsequence converging to the point $ x $ . So we prove
  \[
      \neg B \Rightarrow \neg A \; \neg C \Rightarrow \neg A
  \]

  Start with $ \neg B \Rightarrow \neg A $. Assume that K is not closed. That means that $ \exists \left( x_n \right) \in  K $ such that $ \lim_{n \to \infty} x_n = x \notin K $. In english, this means that there is a boundary point $ x \notin K $. Take any $ \left( x_n \right) \in K  $ converging to $ x $. Then $ \lim_{n \to \infty} x_n = x $. And since $ x_n $ converges to a point $ x $, then so does all its subsequences. So 
  \[
      \lim_{n   \to \infty} x_{n_k} = x 
  \]
  But since K is not closed, there exists limit points that are not in K. So $ x \notin K $, and thus, there exists sequences in K such that there are no subsequences converging to a point in K. Hence, K cannot be compact if K is not closed


  $ \neg C \Rightarrow  \neg A $ 
  Assume now that K is not bounded, that is
  \[
      \forall M \in \mathbb{R} \; \exists a, b \in K \text{ such that } d \left( a, b \right) > M
  \]
  
  Pick $ x_0 \in K $. Since $ \forall n \in \mathbb{N} \; \exists x_n \in  K $ such that $ d \left( x_n, x_0 \right) > n $, then

  \[
      \lim_{n \to \infty}  d \left( x_n, x_0 \right) = \infty
  \]

  We want to show that $ x_n $ has no convergent subsequence. This happens if for $ x \in K $, we have that 

  \[
      \lim_{k \to \infty} d \left( x_{n_k}, x \right) = 0 
  \]

  And from the triangle ineq.
  \[
      d \left( x_{n_k}, x_0 \right) \leq d \left( x_{n_k}, x \right) + d \left( x, x_0 \right) 
  \]
  
  Since $ d \left( x_{n_k}, x \right) \to 0 $ as $ k \to \infty $, we have that $ d \left( x_{n_k}, x_0 \right) \leq d \left( x, x_0 \right) =: C $. But from above, we have that 
  \[
      d \left( x_{n_k}, x_0 \right) > n_k
  \]
  and as $ k \to \infty $ we get that $ d \left( x_{n_k}, x_0 \right) \to \infty $. And hence its a contradiction, since

  \[
      \infty \leq C 
  \]
  is not true.
  
  
  
  
  
  
\end{proofbox}
\end{propositionbox}

\begin{corollarybox}{3.5.5}{}
A subset of $ \mathbb{R}^n $ is compact $ \iff  $ it is closed and bounded

\begin{proofbox}{3.5.5}{}
  We have already proved that $ A \Rightarrow B \& C $, must now prove that a closed and bounded subset of $ \mathbb{R}^n $  is compact. 

  Assume A is closed and bounded. From Bolsano Weierstrass theorem, any bounded sequence in $ \mathbb{R}^n $ has a convergent subsequence. This means if we take $ \left( x_n \right) \in A $ , then

  \[
      \lim_{k \to \infty} x_{n_k} = x
  \]
  This can be viewed as yet another sequence. Define $ \left( y_k \right) = \left( x_{n_k} \right)  $. Then since A is closed, $ x \in K $, and we have found that if we pick an arbitrary sequence in $ A \subseteq \mathbb{R}^n $, then we get that $ \exists \left( x_{n_k} \right)  $ such that $ d \left( y_k, x \right) \to 0 $ as $ k \to \infty $ 
  
\end{proofbox}
\end{corollarybox}

\begin{lemmabox}{3.5.6}{}
Assume that $ \left\{ x_n \right\}  $ is a caucht sequence in a metric space. If there is a subsequence $ \left\{ x_{n_k} \right\}  $ converging to a point $ a $, then the original sequence $ \left\{ x_n \right\}  $ also converges to $ a $.

\begin{proofbox}{3.5.6}{}
  Assume $ \left( x_n \right) \in \left( X, d \right)   $ is cauchy, and that $ \left( x_n \right)  $ has a subsequence $ \left( x_{n_k} \right)  $ such that $ d \left( x_{n_k}, a \right) \to 0 $  as $ k \to \infty $. Since $ \left( x_n \right)  $  is cauchy, we have that $ \forall \epsilon > 0 \; \exists N \in \mathbb{N} $ such that $ d \left( x_n, x_m \right) < \epsilon  \; \forall n, m \geq N $. Since $ \left( x_{n_k} \right)  $   converges to $ a $, there is a $ K \in \mathbb{N} $ such that $ n_K \geq N $ and $ d \left( x_{n_k}, a \right) < \epsilon  $. Then
  \[
      d \left( x_n, a \right) \leq d \left( x_n, x_{n_K} \right) + d \left( x_{n_K}, a \right) < \epsilon + \epsilon = 2 \epsilon 
  \]
  
\end{proofbox}
\end{lemmabox}

\begin{propositionbox}{3.5.7}{}
Every comapct metric space is complete

\begin{proofbox}{3.5.7}{}
  Assume that K is compact. Then $ \forall \left( x_n \right) \; \exists \left( x_{n_k} \right)  $ such that $ \left( x_{n_k} \right)  $ converges to $ x \in K $. K is complete is all cauchy sequences converge.

  Assume now that $ \left( x_n \right)  $ is a cauchy sequence in compact metric space $ \left( X, d \right)  $. Since X is compact, $ \left( x_n \right)  $ has convergent subsequences. And from the lemma above, if a cauchy sequence has convergent subsequences, then $ \left( x_n \right)  $ also converges to $ x \in X $. Hence all cauchy sequences converge, and X is complete as well
\end{proofbox}
\end{propositionbox}


\begin{propositionbox}{3.5.8}{}
A closed subset $ F $ of a compact set $ K $ is compact

\begin{proofbox}{3.5.8}{}
  Assume F is closed. Then all convergent subsequences in F converges to $ a \in F $. Assume $ \left( x_n \right)  $ is a sequence in F. Since K is compact, every sequence $ \left( x_n \right)  $ has convergent subsequences $ \left( x_{n_k} \right)  $ to $ a \in L $. 

  But since F is closed, any sequence must stay within the closed set. which means that any subsequence must also stay within F. Then $ a \in F $, and thus, $ \left( x_{n_k} \right)  $  converges to $ a \in F $, which means that an arbitrary sequence $ \left( x_n \right) \in F  $ has a convergent subsequence $ \left( x_{n_k} \right)  $ converging to $ a \in F $, which means that F is also compact. 
\end{proofbox}
\end{propositionbox}

\begin{propositionbox}{3.5.9}{}
Assume that $ f : X \to Y $ is a continuous function between two metric spaces. If $ K \subseteq X $ is compact, then $ f \left( K \right)  $ is a compact subset of $ Y $ 

\begin{proofbox}{3.5.9}{}
  Let $ \left( y_n \right)  $ be a sequence in $ f \left( K \right)$. We have that $ K \subseteq X $  is compact. Since $ f(x) \in f \left( K \right)  $, we have for each $ n \in \mathbb{N} \; f \left( x_n \right) = y_n$. Since K is compact, $ \left( x_n \right)  $ has a converging subsequence $ \left( x_{n_k} \right)  $  to $ x \in K $. But since f is a continuous function, then taking the function of a convergent sequence will preserve convergence in the new metric space. Thus, $ f \left( x_{n_k} \right)  $ converges to $ f \left( x \right)  $

  \[
      \left( y_{n_k} \right) = \left( f \left( x_{n_k} \right)  \right) 
  \]

  and thus, there is a subsequence of $ \left( y_n \right)  $ converging to $ y = f(x) $, and Y is compact as a result
  
\end{proofbox}
\end{propositionbox}


\begin{theorembox}{3.5.10 Extreme value theorem}{}
Assume $ K $ is nonempty compact subset of $ \left( X, d \right)  $ and that $ f : K \to \mathbb{R} $ is a continuous function. Then f has a maximum and minimum points in $ K $, i.e. $ \exists  c, d \in K $ such that
\[
    f \left( d \right) \leq f \left( x \right) \leq f \left( c \right) \; \forall x \in K
\]


\begin{proofbox}{3.5.10}{}
  Since K is compact, its closed and bounded. Define now
  \[
      \alpha_{  }^{} := \inf \left\{ x \in K : f(x) \right\} \; \beta_{}^{} := \sup \left\{ x \in X : f(x) \right\} 
  \]
  then 
  \[
      \alpha_{}^{} \leq f(x) \leq \beta_{}^{} \; \forall x \in  K
  \]
  
\end{proofbox}
\end{theorembox}

\begin{definitionbox}{3.5.11}{}
A subset $ A $ of $ X $ is called totally bounded if for each $ \epsilon > 0 $ there is finite $ B \left( a_1, \epsilon  \right), V \left( a_2, \epsilon  \right), ..., B \left( a_n, \epsilon  \right)    $ of balls with centers in $ A $ and radius $ \epsilon  $ that cover $ A $
\[
    A \subseteq B \left( a_1, \epsilon  \right) \cup V \left( a_2, \epsilon  \right) \cup ... \cup B \left( a_n, \epsilon  \right)
\]
\end{definitionbox}


\begin{propositionbox}{3.5.12}{}
Let $ K $ be a compact subset of $ X $. Then $ K $ is totally bounded

\begin{proofbox}{3.5.12}{}
  Assume that K is not totally bounded. This means that $ \exists \epsilon > 0 $ such that no finite collection of $ \epsilon  $- balls cover K. We shall make a sequence $ \left( x_n \right)  $ that has no convergent subsequence.

  Pick $ x_1 \in K $, then $ B \left( x_1, \epsilon  \right)  $ does not cover K. Choose $ x_2 \in  K \setminus B \left( x_1, \epsilon  \right)  $, then $ B \left( x_1, \epsilon  \right) \cup B \left( x_2, \epsilon  \right)   $ does not cover K. So we pick $ x_3 \in K \setminus \left( B \left( x_1, \epsilon \right), B \left( x_2, \epsilon  \right)   \right)  $, and get that $ B \left( x_1, \epsilon  \right) \cup B \left( x_2, \epsilon  \right) \cup B \left( x_3, \epsilon  \right)    $ also does not cover K.


  Go down and get that $ x_n \in K \setminus \left( B \left( x_1, \epsilon   \right) \cup ... \cup B \left( x_{n-1}, \epsilon  \right)   \right)  $. Now we have a sequence $ \left( x_n \right)  $ such that $ d \left( x_n, x_m \right) \geq \epsilon \; \forall n, m \in \mathbb{N} $ and $ n \neq m $. Hence $ \left( x_n \right)  $ has no convergent subsequence, and A is not compact. 
\end{proofbox}
\end{propositionbox}

\begin{theorembox}{3.5.13}{}
A subset $ A $ of complete metric space $ X $ is compact $ \iff $ $ A  $ is closed and totally bounded

\begin{proofbox}{3.5.12}{}
  To be continued
\end{proofbox}
\end{theorembox}


\begin{oppgavebox}{3.5.5 }{}
Let $ \left( X, d \right)  $ be a metric space. For a subset $ A \subset X $, let $ \partial A $ denote the boundary of A. $ \bar{A} = A \cup \partial A $ 

\begin{itemize}
  \item a) A subset A of X is called precompact $ \bar{A} $ is compact. Show that A is precompact if and only if all sequences in A have convergent subsequences
  
  \begin{proofbox}{3.5.5 a)}{}

    \begin{itemize}
      \item 1: A is precompact
      \item 2: all sequences in A have convergent subsequences.
    \end{itemize}

    We start with $ 1 \Rightarrow 2 $ 
    Assume that A is precompact. That is $ \bar{A} $ is compact. Want to show that all sequences in A have convergent subsequences. 

    Take an arbitrary sequence in A, $ \left( x_n \right)  $. Then $ \left( x_n \right) \subset \bar{A} $, which is compact. This means that $ \left( x_n \right)  $ has a convergent subsequence $ \left( x_{n_k} \right)  $ to $ x \in  \bar{A} $. Hence all sequences in A has a convergent subsequence


    Now we show the other way $ 2 \Rightarrow 1 $

    Assume that all sequences in A have convergent subsequences. Let $ \left( x_n \right) \subset A $. Then $ \left( x_{n_k} \right)  $ converges to $ x \in  X $. But since $ A \subseteq \bar{A} $, then $ x \in \bar{A} $, and hence, since $ \bar{A} $ is closed by definition (the set $\bar{A}$ contains its boundary points) $\left( x_{n_k} \right)  $ must converge to $ x \in \bar{A} $. Thus, $ \bar{A} $ is compact $\Rightarrow$ A is precompact 
  \end{proofbox}
  \item b) Show that a subset of $ \mathbb{R}^n $ is precompact if and only if its bounded
  
  \begin{proofbox}{3.5.5 b)}{}
    We have that 
    \begin{itemize}
      \item A: $ A \subseteq \mathbb{R}^n $ is precompact
      \item B: $ A \subseteq \mathbb{R}^n $ is bounded
    \end{itemize}


    $ A \Rightarrow B $ 

    Assume that $ A \subseteq \mathbb{R}^n$ is precompact.  Since $ \bar{A} \subseteq \mathbb{R}^n $ is compact, then this implies that $ \bar{A} $ is closed and bounded. We already know that $ \bar{A} $ is closed, but now its bounded as well. So $ \exists M \in \mathbb{R} $ such that $ d \left( a, b \right) \leq M \; \forall a, b \in  \bar{A} $. Furthermore, since $ A \subseteq \bar{A} $, then M also encloses $ A $, and we get that $ A $ is also bounded. 



    $ B \Rightarrow A $
    
    Assume that $ A \subseteq \mathbb{R}^n $ is bounded. Then we know that $ \bar{A} = A \cup \partial A $ is closed and bounded. In $ \mathbb{R}^n $, this implies that $ \bar{A} $ is also compact.  So $ \bar{A} $ is compact $ \iff $ A is precompact
    
    And the proof is finished

  \end{proofbox}
\end{itemize}
\end{oppgavebox}


\begin{oppgavebox}{3.5.6}{}
Assume $ \left( X, d \right)  $ is a metric space and $ f : X \rightarrow \left[ 0, \infty \right)  $ is a continuous function. Assume $ \forall  \epsilon > 0 \; \exists  K_{ \epsilon } \subseteq K $ such that $ f(x) < \epsilon  $ when $ x \notin K_{ \epsilon } $. Show that f has a maximum.

\begin{proofbox}{3.5.6}{}
  We have that $ X = K_{ \epsilon } \cup \left( K_{ \epsilon } \right)^c  $. We assume that $ \forall  \epsilon > 0 \; \exists K_{ \epsilon } \subseteq X $ such that $ f(x) < \epsilon  $ when $ x \in \left( K_{ \epsilon }  \right)^c$  .


  When $ x \in K_{ \epsilon } $, then, since $ K_{ \epsilon } $ is compact and f is continuous, we get that there exists a $ d \in K_{ \epsilon } $ such that $ f(x) \leq f(d) \; \forall x \in K_{ \epsilon } $. But we are asked to find the maximum of f in the whole metric space x. So when $ x \notin K_{ \epsilon } $, then $ f(x) < \epsilon  $ and $ \epsilon   $ becomes the upper bound for f. Thus, we can define 

  \[
      \sup_{ x \in X} f(x) = max \left( f(d), \epsilon \right) 
  \]
  What is bigger between $ f(d)  $  and $ \epsilon  $?

  Pick $ \epsilon _1 $ and define
  \[
      \epsilon _n = e^{-n} \epsilon _1
  \]
  Then 
  \[
      \lim_{n \to \infty} \epsilon _n = 0
  \]
  
  Thus, $ \exists N \in \mathbb{N} $ such that $  d \left( \epsilon _n, 0 \right) < d f(d) \; \forall n \geq N $. So we have found that

  \[
      \sup_{x \in X} = f(d) \; \forall n \geq N
  \]
\end{proofbox}
\end{oppgavebox}

\begin{oppgavebox}{3.5.7}{}
Let $ \left( X, d \right)  $ be a compact metric space, and assume that $ f: X \rightarrow \mathbb{R} $ is continuous when we give $ \mathbb{R} $ the usual metric $ d \left( x, y \right) = \left| x - y \right|  $. Show that if $ f(x) > 0 \; \forall  x \in X $, then $ \exists a \in \mathbb{R} \; f(x) > a \; \forall x \in X $.

\begin{proofbox}{3.5.7}{}
  Since $ X $ is compact and f is continuous, then there exists a $ c \in X $ such that $ f(c) \geq f(x) \; \forall x \in X $. Since $ f(x) > 0 \; \forall  x \in X $  then also $ f(c) > 0 $. Hence there exists $ a \in \mathbb{R} $ such that
  \[
      f(c) = a \leq f(x) \; \forall x \in X
  \]
\end{proofbox}
\end{oppgavebox}

\begin{oppgavebox}{3.5.8}{}
Assume that $ f : X \rightarrow Y $ is continuous between metric spaces, and let $ K $ be a compact subset of Y. Show that $ f^{-1} (K) $ is closed. Find an examole which shows that $ f^{-1}(K) $ need not be compact

\begin{proofbox}{3.5.8}{}
  Since $ K \subseteq Y $ is compact, K is closed and bounded. So all sequences $ \left( y_n \right) \in K $ have convergent subsequences meaning that $ \left( y_{n_k} \right)  $ converges to $ y \in K $. Since f is continuous, we have that $ \forall n \in \mathbb{N} $, we have $ f^{-1}(y_n) = x_n $. This means that 
  \[
      f^{-1} \left( y_{n_k} \right) = x_{n_k}
  \]
  Now taking the limits, we get that $ \lim_{k \to \infty} f^{-1} \left( y_{n_k} \right) = f^{-1} \left( \lim_{k \to \infty} y_{n_k} \right) = f^{-1} (y) = x \in f^{-1} \left( K \right) $. Since this counts for all $ \left( y_n \right)  $  in $ K $, then $ f^{-1} (K) $ is closed.   
\end{proofbox}
\end{oppgavebox}

\begin{oppgavebox}{3.5.9}{}
Show that a totally bounded subset of a metric space is always bounded. Find and example of a bounded set in a metric space that is not totally bounded.

\begin{proofbox}{3.5.9}{}
  Assume $ A \subseteq X $ is totally bounded, then if $ \exists \epsilon > 0 $ such that for a finite collection of balls,  $ A \subseteq B \left( a_1, \epsilon  \right) \cup ... \cup B \left( a_n, \epsilon  \right)   $. 

  Choose $ \epsilon > 0 $. We have that $ A \subseteq \cup_{i=1}^n B \left( a_i, \epsilon  \right)  $. Pick an arbitrary $ x_0 \in X $, then define $ R := max_{i \in \left[ 1, n \right] } d(a_i, x_0 ) + 1$, thus $ d \left( a_i, x_0 \right) < R \; \forall i \in \left[ 1, n \right]  $ . Pick $ x \in A $ such that $ x \in B \left( a_i, \epsilon  \right)  $ , then $ d(a_i, x) < \epsilon  $. we get
  
  \[
      d \left( x, x_0 \right) \leq d \left( x, a_i \right) + d \left( a_i, x_0 \right) < \epsilon + R
  \]

  hence $ x \in B \left( x_0, R \right)  $ 
  
\end{proofbox}
\end{oppgavebox}


\begin{oppgavebox}{3.5.10}{}
Let $ d \left( x, y \right) = \left| x - y \right|  $ be the usual metric on $ \mathbb{Q} $. Let $ A = \left\{ q \in \mathbb{Q} : 0 \leq q \leq 2 \right\}  $. Show that A is closed and totally bounded subset of $ \mathbb{Q} $, but not compact.

\begin{proofbox}{3.5.10}{}
  A is not compact, since the sequence converging to $ \sqrt{ 2 } $ does not lie in $ \mathbb{Q} $, and hence not in $ A $. But its closed, because $ A^c = \left\{ q \in \mathbb{Q} : \left( - \infty, 0 \right) \cup \left( 2, \infty \right)   \right\}  $ is open. Its totally bounded since 
\end{proofbox}
\end{oppgavebox}

\begin{oppgavebox}{3.5.11}{}
A metric space $ \left( X, d \right)  $ is locally compact if for each $ a \in X $ is an $ r > 0 $ such that the closed ball $ \bar{B} \left( a, r \right) = \left\{ x \in X : d \left( a, x \right) \leq r \right\}   $ is compact.
\begin{itemize}
  \item a) Show that $ \mathbb{R}^n $ is locally compact.
  \begin{proofbox}{3.5.11 a)}{}
    Take any $ x \in \mathbb{R}^n $, then the closed ball around x with radius r is $ \bar{B} \left( x, r \right)  $. Show that $ \bar{B} \left( x, r \right)  $ is compact. 

    Since $ \bar{B} \left( x, r \right) \subseteq \mathbb{R}^n $, we have to show that $ \bar{B} \left( x, r \right)  $ is compact. This happens if its closed and bounded. Its already closed by definition, so we only have to show its bounded.
    

    BOUNDEDNESS:
    $ \exists M \in \mathbb{R} $ such that $ d \left( a, b \right) \leq M \; \forall a, b \in \bar{B} \left( x, r \right)  $. Since $ d \left( x, y \right) \leq r \; \forall y \in \bar{B} \left( x, r \right)  $, we have a bounded set. 
  \end{proofbox}

  \item b) Show that if $ X = \mathbb{R} \setminus \left\{ 0 \right\}  $ and $ d : X \rightarrow \mathbb{R} $ is defined as $ d \left( x, y \right) := \left| x-y \right|  $, then $ \left( X, d \right)  $ is locally compact, but not complete.
  
  \begin{proofbox}{3.5.11 b)}{}
    We know that its not complete since the sequence $ x_n = \frac{ 1 }{ n } $ converges to $ 0 \notin \mathbb{R} \setminus \left\{ 0 \right\}  $. Its locally compact if $ \forall x \in \mathbb{R} \setminus \left\{ 0 \right\}  \; \exists \epsilon  > 0 $ such that $ \bar{B} \left( x, \epsilon  \right)  $ is compact. This happens if its closed and bounded. Its already closed from definition, it contains its boudnary points.... Its bounded if $ \exists M \in \mathbb{R} $ such that $ d \left( x,y \right) \leq M \; \forall y, x \in \bar{B} \left( x, eps \right) $  . So we have that $ \forall y \in \bar{B} \left( x, \epsilon  \right) \Rightarrow d \left( x, y \right) = \left| x-y \right| \leq \epsilon =: M$. Hence we have that its closes
  \end{proofbox}
\end{itemize}
\end{oppgavebox}

\begin{oppgavebox}{3.5.15}{}
Assume C and K are disjoint compact subset of a metric space $ \left( X, d \right)  $ and define
\[
    a = \inf \left\{ d \left( x, y \right) : x \in C, y \in K \right\} 
\]

Show that $ a > 0 $ and that $ \exists x_0 \in C, y_0 \in K $  such that $ d \left( x_0, y_0 \right) =a $

\begin{proofbox}{3.5.15}{}
  Since $ C \subseteq X $ and $ K \subseteq X $, then C and K inherits the metric, and $ \left( C, d_C \right)  $ and $ \left( K, d_K \right)  $ are two new metric spaces in $ \left( X, d \right)  $ 
  Construct a new metric space $ c \times K $ with the metric $ d_{\times} \left( \left( x, y \right) , \left( x', y' \right)  \right) = d_C \left( x, x' \right) + d \left( y, y' \right)   $. Since the distance function $ d_C : C \rightarrow \mathbb{R} $ and $ d_K : K \rightarrow \mathbb{R} $  are continuous, then $ f : C \times K \rightarrow \mathbb{R} $ defined by $ f(x,y) = d_{C \times K}(x,y) $, then f is a continuous function. To show that $ a > 0 $, assume that $ a = 0 $. This means that $ d(x,y) = 0 \Rightarrow x,y \in C \cap K = \emptyset $. So the statement that $ d \left( x, y \right) = 0 $ is false. Therefore, $ a > 0 $ (because $ d \left( x, y \right) \geq 0 $)    
  
  
  
  The product of two metric spaces are also compact, and hence, f attains a minimum
  \[
      \exists x_0, y_0 \in C \times K \text{ such that } f \left( x_0, y_0 \right) = min_{x, y \in C \times K} f(x, y)
  \]
   
  We now just let $ a = f \left( x_0, y_0 \right)  $ and get that $ d \left( x_0, y_0 \right) = a $ 

\end{proofbox}


\begin{proofbox}{Proof that the product of two compact metric spaces are also compact}{}
  Let $ \left( X, d_X \right)  $ and $ \left( Y, d_Y \right)  $ be two compact metric spaces. We have that  
  \[
      d_X : X \rightarrow \mathbb{R}
  \]
  \[
      d_Y : Y \rightarrow \mathbb{R}
  \]
  The product of the two metric spaces make up $ \left( X \times Y, d_{X \times Y} \right)  $. Is this new metric space compact?

  We have that since $ \left( X,d_X \right)  $ and $ \left( Y, d_Y \right)  $ are compact, then
  \[
      \left( x_n, y_n \right) 
  \]
  have convergent subsequences
  \[
      \left( x_{n_k}, y_n \right) \to \left( x, y_n \right) 
  \]
  and
  \[
      \left( x_n, y_{n_k} \right) \to \left( x_n, y \right) 
  \]

  hence
  \[
      d_{X \times Y} \left( \left( x_{n_k}, y_{n_k} \right), \left( x, y \right)  \right)  \to d_X \left( x_{n_k}, x \right) + d_Y \left( y_{n_k}, y \right) = 0
  \]
  So $\left( x_{n_k}, y_{n_k} \right) \to \left( x, y \right) $ as $ k \to \infty $ 
  
\end{proofbox}
\end{oppgavebox}

\begin{oppgavebox}{A metric space}{}
Define a metric space as rigorously as possible
\begin{proofbox}{}{}
  A metric space $ \left( X, d_X \right)  $ is a non-empty set that satisfies
  \begin{itemize}
    \item $ d \left( x, y \right) \geq 0  $ and $ d \left( x, y \right) = 0 \iff x =y \; \forall x, y \in X $
    \item $ d \left( x,y \right) = d \left( y, x \right) \; \forall x, y \in X $
    \item $ d \left( x, y \right) \leq d \left( x, z \right) + d \left( z, y \right) \; \forall x, y, z \in X  $   
  \end{itemize}
\end{proofbox}
\end{oppgavebox}

\begin{oppgavebox}{A subset of metric space is bounded}{}
\begin{proofbox}{}{}
  $ A \subseteq \left( X, d \right)  $ is boudned if $ \exists M \in \mathbb{R} $ s.t. $ d \left( a, b \right) \leq M \; \forall a, b \in A $  
\end{proofbox}
\end{oppgavebox}

\begin{oppgavebox}{An open subset of a metric space}{}
\begin{proofbox}{}{}
  An open subset $ A \subseteq X $ of a metric space $ \left( X, d \right)  $ is open if and only if $ \forall x \in A \; \exists \epsilon > 0 $ such that $ B \left( x, \epsilon \right) = \left\{ y \in A : d \left( x, y \right) < \epsilon  \right\} \subseteq A  $ 
\end{proofbox}
\end{oppgavebox}

\begin{oppgavebox}{A closed subset of a metric space}{}
Write as rigorously as possible the definition

\begin{proofbox}{}{}
  A closed subset $ A \subseteq X $ of metric space $ \left( X, d \right)  $ is defined as 
  
  $ \forall \left( x_n \right) \in A $ such that $ \left( x_n \right) \to x \in X \Rightarrow x \in A $
\end{proofbox}

\begin{proofbox}{Alternative proof}{}
  A contains all the limit points of any convergent sequence. That is, if $ x \in X $ satisfies
  \[
      \forall \epsilon > 0 \; B \left( x, \epsilon \right) \cap \left( A \setminus \left\{ x \right\}  \right)  = \emptyset
  \]
  then $ x \in A $ 
  
\end{proofbox}

\begin{proofbox}{Alternative}{}
  $ A \subseteq X $ is closed $ \iff \; A^c = X \setminus A $ is open. That is, $ \forall x \in X \setminus A \; \exists  \epsilon > 0 $ such that $ B \left( x, \epsilon \right) \subseteq X \setminus A  $  
\end{proofbox}
\end{oppgavebox}

\begin{oppgavebox}{Convergence of a sequence in a metric space}{}
\begin{proofbox}{}{}
  $ \left( x_n \right) \in X $ converges to x if $ \forall \epsilon > 0 \; \exists  N \in \mathbb{N} $  such that $ d \left( x_n, x \right) < \epsilon \; \forall n \geq N $ 
\end{proofbox}

\begin{proofbox}{}{}
  $ \left( x_n \right) \to x $ as $ n \to \infty $ 
\end{proofbox}
\end{oppgavebox}

\begin{oppgavebox}{A function f from one metric space to another is continuous at a point x}{}
\begin{proofbox}{}{}
  Let $ \left( X, d_X \right)  $ and $ \left( Y, d_Y \right)  $ be two metric spaces. f is continuous at $ x \in X $ if $ \forall \epsilon > 0 \; \exists \delta > 0  $ such that $ \forall y \in X \; d \left( x, y \right) < \delta \Rightarrow d \left( f(x), f(y) \right) < \epsilon $ 
\end{proofbox}

\begin{proofbox}{}{}
  $ \forall \left( x_n \right) \to x \in X \; f \left( x_n \right) \to f(x) $
\end{proofbox}
\end{oppgavebox}

\begin{oppgavebox}{A sequence in a metric space is cauchy}{}
\begin{proofbox}{}{}
  Let $ \left( X, d \right)  $ be a metric space, $ \forall \epsilon > 0 \; \exists N \in \mathbb{N} $ such that $ d \left( x_n, x_m \right) < \epsilon \; \forall n, m  \geq N $ 
\end{proofbox}
\end{oppgavebox}

\begin{oppgavebox}{A complete metric space}{}
\begin{proofbox}{}{}
  A metric space $ \left( X,d \right)  $ is complete if all cauchy sequences converge. That is, let $ \left( x_n \right)  $ be a cauchy sequence in $ \left( X, d \right)  $, then   $ \forall \epsilon > 0 \; \exists  N \in \mathbb{N} $ such that $ d \left( x_n, x_m \right) < \epsilon \; \forall n, m \geq N $ 
\end{proofbox}
\end{oppgavebox}

\end{document}
